\documentclass[11pt,letterpaper]{article}
\usepackage[margin=1.04in,top=.85in,bottom=.9in]{geometry}
\usepackage{amsmath,amssymb,amsthm,mathtools}
\usepackage{booktabs,tabularx,array,xcolor,enumitem,fancyhdr}
\usepackage[T1]{fontenc}
\usepackage{lmodern}
\usepackage[breaklinks=true,colorlinks=true]{hyperref}
\definecolor{ink}{HTML}{17354D}\definecolor{sea}{HTML}{176B74}
\hypersetup{linkcolor=sea,urlcolor=sea,citecolor=sea,pdftitle={Quantum mechanics in intuitionistic primitive recursive arithmetic},pdfauthor={OpenAI}}
\usepackage{titlesec}
\titleformat{\section}{\Large\bfseries\color{ink}}{\thesection}{.7em}{}
\titleformat{\subsection}{\large\bfseries\color{ink}}{\thesubsection}{.65em}{}
\setlist{topsep=.35em,itemsep=.15em}
\setlength{\parindent}{1.1em}\setlength{\parskip}{.15em}
\setlength{\emergencystretch}{2.2em}
\setlength{\headheight}{14pt}
\pagestyle{fancy}\fancyhf{}\fancyhead[L]{\small\color{ink}Quantum mechanics in iPRA}\fancyhead[R]{\small\color{ink}Mathematical foundations}\fancyfoot[C]{\small\thepage}
\renewcommand{\headrulewidth}{.25pt}
\newtheorem{theorem}{Theorem}[section]\newtheorem{lemma}[theorem]{Lemma}
\newtheorem{proposition}[theorem]{Proposition}\newtheorem{corollary}[theorem]{Corollary}
\theoremstyle{definition}\newtheorem{definition}[theorem]{Definition}\newtheorem{example}[theorem]{Example}\newtheorem{remark}[theorem]{Remark}
\setcounter{tocdepth}{1}\allowdisplaybreaks[2]\raggedbottom
\begin{document}
\thispagestyle{plain}
\noindent{\small\sffamily\bfseries\color{sea} CONSTRUCTIVE ANALYSIS\quad /\quad MATHEMATICAL PHYSICS}\par\vspace{1em}
\noindent{\fontsize{29}{33}\selectfont\bfseries\color{ink} Quantum mechanics\par}
\noindent{\fontsize{18}{21}\selectfont\bfseries\color{ink}\mbox{in intuitionistic primitive recursive arithmetic}\par}
\vspace{.7em}\noindent{\large Quantitative proofs for states, measurements and dynamics}\par
\vspace{.6em}\noindent{\small 20 November 2026}\par\vspace{1em}
\noindent\textbf{Abstract.} An experiment specified by preparation, evolution and an effect should determine its probability to any requested precision. We give a first-order arithmetic meaning to this statement for nonrelativistic quantum mechanics. Finite-dimensional states and channels admit positive square roots, purifications and Kraus representations with primitive-recursive precision schedules, including at singular matrices. For a separable Hilbert space, basis tails determine the information needed to evaluate operators, traces and countable experiments. We construct the continuous functional calculus and give a quantitative passage from resolvents of an unbounded Hamiltonian to its unitary group and its exact generator domain. The oscillator and the free particle exhibit discrete and continuous energy measurements. An undecidability argument already in dimension one shows why arbitrary sharp spectral cuts require additional data. The archive contains exact examples and Lean object-language derivations for selected arithmetic lemmas. The general analytic proofs are given on paper, with a formalization plan.
\vspace{.25em}\begingroup\setlength{\parskip}{0pt}\tableofcontents\endgroup\newpage
\section{The arithmetic meaning of a quantum prediction}\label{sec:arithmetic}
Consider a spin prepared in amplitudes $(z,w)$, sent through a Hadamard gate and then measured in the original basis. The first outcome has probability $|z+w|^2/2$. It matters whether a device adds the amplitudes before squaring. For $z=w=1/\sqrt2$ the first outcome is certain, although the incoherent mixture with the same initial diagonal has first-outcome probability $1/2$. The identity
\[
 |z+w|^2+|z-w|^2=2(|z|^2+|w|^2)
\]
certifies normalization without evaluating a square root. Clearing denominators turns this conservation law into an integer polynomial identity. The same principle of finite observation works for a wave packet with infinitely many momentum values, once a quantitative rule supplies its neglected norm. The significant mathematical work is to construct such rules and show that predictions depend only on the represented objects.

\subsection{Theory, codes and logic}
We use the ordinary iPRA of ProveIt \cite{proveit}. Its domain consists of natural numbers. The language has zero, successor, and a symbol for each primitive recursive (PR) definition. Its axioms give the defining equations, successor injectivity and separation from zero, and induction on quantifier-free formulas with free numerical parameters. The logic is intuitionistic first-order predicate logic with equality. Parameters, requested precisions, integer dimensions, time bounds and cutoff indices are all numerical. ``For every presented Hilbert vector'' is a metatheoretic schema: fix the PR symbols for its samplers, introduce any numerical parameters as free variables, and assert an implication from the stated first-order laws. Every explicit numerical family with one joint PR sampler gives a single parameterized instance. For an operation described as a supplied evaluator on effects or names, require that applying it to each joint PR family used in the proof gives a joint PR family, with the stipulated laws. An evaluator on arbitrary natural-number codes for PR programs would be a different and stronger assumption. A law supplied as a hypothesis remains a hypothesis until its instance has a derivation.

Here is one concrete choice of coding. Represent a signed integer by $(p,n)$ with value $p-n$, a rational by $(p,n,d)$ with value $(p-n)/(d+1)$, and a complex rational by two rationals. Equality and order reduce to natural-number comparisons by cross multiplication. A length and a nested pair of entries encode a finite list. For example $\pi(u,v)=(u+v)^2+u$ has projections found by two bounded loops up to the input code; give unused codes the default value zero. Reading a list, scanning a matrix, summing and multiplying by a length bound, concatenation and building a code as a loop accumulator are PR. A maximum, a least index within a stated bound, and a bounded universal test are PR (fold the conjunction of the individual Boolean answers). The same operations implement arbitrary rectangular shapes, including empty environments when an empty sum is appropriate. Positive denominators can always be retained by products, without reduction to lowest terms. Dimensions of physical state spaces are positive.

Write $\delta_k=2^{-k}$. A real or complex scalar is a fixed PR sampler of rational values $a_k$ satisfying $|a_k-a_l|\le\delta_k$ for $l\ge k$, with an integer magnitude bound. For complex values use the squared absolute value to check the inequality. Equality means $|a_k-b_k|\le2\delta_k$ for every $k$. The closed order on reals is $a_k\le b_k+2\delta_k$ for every $k$. An apartness or positivity witness gives a specified positive rational lower bound, possibly after computing more digits. Sample values that are certified to within $\delta_{k+2}$ of the intended answer at stage $k$ define a regular name. Section~\ref{sec:infinitehilbert} proves the corresponding results for vector names. An operator estimate $\|T\|\le b$ denotes an inequality on all its finite-vector tests. It never requires forming a real supremum.

Scalar addition and multiplication use the triangle inequality and the estimate
$|ab-a'b'|\le |a||b-b'|+|a-a'||b'|$. To compute a positive square root on $[0,B]$, approximate the input to $\delta_{2k+6}$, clamp the rational to $[0,B]$, and bisect the interval $[0,1+B]$ for its square root to width $\delta_{k+3}$. This costs at most $k+4+\lceil B\rceil$ steps. The bound $|\sqrt u-\sqrt v|^2\le|u-v|$ follows by factoring the difference of squares. This reasoning may first define square roots of rationals by the bisection itself, and then extend to names. Division by a real of magnitude at least $b>0$ is PR in its name and the rational witness, since input error $\eta\le b/2$ gives reciprocal error at most $2\eta/b^2$. For complex division use the conjugate divided by the squared modulus. These rules also yield absolute values and rational upper bounds for subsequent calculations. No sign decision on a general named real is used.

\subsection{From estimates to first-order arguments}
There are three recurring moves in the proofs.
\begin{enumerate}
\item A closed inequality at the limit is proved to arbitrary rational slack. If $r,s$ are rational observations and $r\le s+C2^{-m}$ for every $m$, a violation $r-s=u/v>0$ gives a contradictory finite test at $m=1+\lceil Cv\rceil$ for integers $u,v\ge1$. When $r,s$ themselves are named, first fix their observation stage and include its prescribed errors. Thus removal of slack uses only rational decidability and bounded integer arithmetic.
\item A limit is constructed only with a joint PR sampler and a PR Cauchy rate. Given the rate $m(k)$, take the $(k+3)$rd sample at index $m(k+3)$ (use running maxima for the rate). The three errors, from the first sample, the Cauchy comparison and the second sample, establish regularity. The same argument proves a uniform operator estimate by testing one finite input vector at a time, scaling the tolerance with its known bound.
\item For a loop of length $n$, fix its requested output tolerance before induction. Rounding each addition to error $\eta$ contributes at most $n\eta$. If a step amplifies error by $L\ge1$, the accumulated error with per-step residual $\eta$ is at most $nL^n\eta$ plus $L^n$ times the initial error. Prove the finite prefix inequality by quantifier-free induction on the actual rational arrays with these fixed parameters. For nonuniform loops replace the powers by products of the supplied integer step bounds. When a hypothesis is a law at every precision, instantiate it at the finitely many precisions needed by the computation. Bounded conjunctions can be folded into Boolean flags. Quantify the output tolerance after this finite proof. In particular we do not induct on ``the limit exists'' or ``this circuit is unitary at all precisions.''
\end{enumerate}
Two details are important for uniformity. A backward pass through a finite rational expression computes all needed precision requests once forward bounds on intermediate magnitudes have been collected. When a variable-length product uses a fixed supplied strong operator $A$, evaluate it by iterating on \emph{finite rational vectors}: approximate $Aq$ to a specified local error from $q$, save the finite result as the next accumulator, and repeat. For example a power of degree $n$ with $\|A\|\le b$ takes per-step error at most $\delta_{k+3}/((n+1)(1+b)^n)$ for absolute output error below $\delta_{k+2}$ given an equally controlled input. This is ordinary numerical primitive recursion. Iteration of a procedure that builds and later evaluates arbitrary function syntax would require a different argument. There is no PR evaluator for all PR programs, by diagonalization. Our schemata compile fixed input programs by substitution. A variable list of gates uses a single joint sampler indexed by the gate number or a fixed finite gate alphabet.

\subsection{What serves as a physical postulate}
Operationally one specifies preparations, admissible effects and an evolution. Their interpretation uses the complex Hilbert-space rules: composition of independent systems is tensor product, isolated reversible time evolution is unitary, and probabilities on compatible outcomes add. These are modeling assumptions about experiments. In finite dimension, assuming a single normalized additive probability assignment on \emph{all effects} already implies the density-matrix trace rule (Theorem~\ref{thm:q-effect}). The choice of Hamiltonian and its parameters is further physical input. We use $\hbar=1$ throughout the dynamics; in physical units the generator is $-iH/\hbar$, with a positive lower bound for $\hbar$ if it is itself named.

An experiment here produces a probability, so its definition does not require a source of random bits. To connect that number with a frequency, take $N$ independent tensor preparations and repeat a two-outcome effect of probability $p$. For the indicator variables $X_j$, tensor factorization gives $\mathbb E X_j=p$ and $\mathbb E X_iX_j=p^2$ for $i\ne j$. Expanding the finite square gives $\mathbb E(\overline X-p)^2=p(1-p)/N\le1/(4N)$. For a rational target $r$ with a certified $|p-r|\le\epsilon/2$, the event $|\overline X-r|\ge\epsilon$ is decidable on its finite outcome strings. On it $(\overline X-p)^2\ge\epsilon^2/4$, hence its total weight is at most $1/(N\epsilon^2)$. An integer $N\ge1/(\alpha\epsilon^2)$ suffices for failure probability $\alpha>0$. This finite-sample assertion also applies when $p$ is an infinite-dimensional Born probability supplied by the later constructions.

\section{States, operations and experiments in finite dimension}
\label{sec:finite-quantum}

Take a positive integer \(d\) and the standard basis \(e_1,\ldots,e_d\).
A vector consists of \(d\) named complex numbers. We use
\(\langle x,y\rangle=x^*y\), linear in the second variable. A rectangular
matrix acts by finite sums. We write \(\|T\|\le b\) for the estimate
\(\|Tx\|_2\le b\|x\|_2\) for all test vectors; no operation taking a
supremum is needed. Rational complex test vectors suffice, by approximation.
The order \(H\ge0\) means \(H=H^*\) and \(x^*Hx\ge0\) on these tests.
An effect is \(0\le E\le I\). A pure vector has \(x^*x=1\). Its state is
the matrix \(xx^*\), which is invariant under multiplication of \(x\) by a
scalar of modulus one.

\paragraph{A precision convention.}
If each entry of an \(m\)-by-\(d\) matrix has complex error at most
\(\eta\), its error on a unit vector is at most \(\sqrt{md}\,\eta\)
by rowwise Cauchy--Schwarz. The integer \(md\) also suffices for positive
dimensions. If \(S,T\) have respective integer norm bounds \(a,b\) and
both approximations have norm error at most \(\eta\le1\), the product
error is at most \((a+b+1)\eta\), by inserting the intermediate product
\(S\widetilde T\). For a sum of \(r\) terms the error is at most
\(r\eta\). A square trace has error at most \(d\eta\) for norm error
\(\eta\), by testing the standard basis. Each guard precision in these
rules can be taken as the requested precision plus the integer loss factor,
since \(2^k\ge k\) for positive integers. In products a bound for the
next exact value is \(ab\), and in sums it is the sum of the bounds.
On a finite arithmetic expression first propagate these bounds forward,
then propagate requested precisions backward. Integer recursion over the
expression implements both passes. All claims below about polynomial
evaluation can use this deliberately coarse rule; it permits coefficients
and matrix sizes to vary with the input. An exact equality of named
quantities always means equality at arbitrarily small rational tolerance.
For example, a telescoping argument invoking a law at each of \(r\) steps
uses that law at rational discrepancy \(2^{-p}/(r+1)\) per step.
With the output precision and \(r\) fixed, the partial-sum error bound
is a rational comparison on finite codes. Quantifier-free induction acts
on that comparison. The universal precision quantifier is introduced after
the induction. This applies also to repeated effect additivity below.

\subsection{Positivity without diagonalization}

\begin{lemma}\label{lem:q-forms}
A positive Hermitian form \(h\) satisfies
\( |h(x,y)|^2\le h(x,x)h(y,y)\).
In particular \(0\le X\le I\) implies \(\|X\|\le1\).
\end{lemma}
\begin{proof}
For positive rational \(\epsilon\), substitute
\(y-xh(x,y)/s\), with \(s=h(x,x)+\epsilon\), in the positivity
inequality. Expansion gives
\((2s-h(x,x))|h(x,y)|^2\le s^2 h(y,y)\).
As \(2s-h(x,x)\ge s>0\), we get \(|h(x,y)|^2\le s h(y,y)\).
Let \(\epsilon=2^{-k}\); the excess on the right is bounded by
\(2^{-k}\) times any integer bound for \(h(y,y)\). For the second
assertion put \(h(x,y)=x^*Xy\). Its absolute value is at most
\(\|x\|_2\|y\|_2\). Taking the first argument to be \(Xy\) yields
\(\|Xy\|_2^2\le\|Xy\|_2\|y\|_2\). If the claimed bound were violated
by a positive rational margin then \(\|Xy\|_2\) would have a positive
lower bound and could be cancelled. This excludes every such margin,
which is precisely the closed order on named reals.
\end{proof}

\begin{theorem}[Positive square root]\label{thm:q-sqrt}
Given \(H\ge0\) with an integer \(L\ge1\) satisfying
\(H\le L^2 I\), one can construct the unique \(S\ge0\) with \(S^2=H\).
The approximation with \(m\ge1\) terms has operator error at most
\(L/\sqrt{m+1}\). Every operator commuting with \(H\) commutes with \(S\).
\end{theorem}
\begin{proof}
Set \(X=I-H/L^2\), a positive contraction. Define rational numbers
\[
 c_0=1,\qquad c_j=4^{-j}{2j\choose j},\qquad
 b_j=c_{j-1}-c_j=c_{j-1}/(2j)\quad(j\ge1),
 \qquad F_m(z)=1-\sum_{j=1}^m b_jz^j.
\]
From \(c_{j+1}/c_j=(2j+1)/(2j+2)\), induction gives
\(c_j^2\le 1/(j+1)\): after squaring the step, the required integer
inequality is \((2j+1)^2(j+2)\le4(j+1)^3\).
All powers of \(X\) are contractions. Even powers are positive by a
square, and \(x^*X^{2s+1}x=(X^s x)^*X(X^s x)\ge0\). Thus
\(c_m I\le F_m(X)\le I\). For \(l\ge m\) the norm of
\(F_l(X)-F_m(X)\) is at most \(c_m-c_l\).

Here is a finite verification of the squared polynomial identity used
in the limit. The coefficient recurrence gives
\((1-z)F_m'=-F_m/2\) through degree \(m-1\).
Therefore \(G=F_m^2\) satisfies \((1-z)G'=-G\) to that order.
With \(G(0)=1\), comparison of consecutive coefficients gives
\(G_1=-1\) and \(G_2=\cdots=G_m=0\).
Every coefficient of \(G\) above degree \(m\) is nonnegative since it
comes from the product of two negative coefficients. Their sum is
\(G(1)=c_m^2\). Hence \(\|F_m(X)^2-(I-X)\|\le c_m^2\).
The limit \(S\) of \(L F_m(X)\) has all the claimed properties except,
as yet, uniqueness. For operator truncation error \(2^{-p-1}\) take
\(m=L^2 2^{2p+2}\); the other half of the budget covers finite arithmetic.

For uniqueness we first record a consequence of the construction.
If positive \(C,D\) commute, the constructed square root \(R\) of \(C\)
commutes with \(D\) by its polynomial formula, so \(CD=RDR\ge0\).
A second positive root \(T\) commutes with \(H=T^2\), hence with \(S\).
Write \(U=S-T\), \(W=S+T\). Then \(UW=WU=0\), and
\(W^2-U^2=4ST\ge0\). Conjugation by \(U\) gives \(-U^4\ge0\);
also \(U^4\ge0\). Evaluating on \(x\) gives \(\|U^2x\|_2^2=0\);
then \(\|Ux\|_2^2=x^*U^2x=0\). This proves uniqueness.
\end{proof}

An integer \(L\) can be obtained from the dimension and coarse entry
bounds. The proof works verbatim for a bounded self-adjoint operator on
a named Hilbert space when those polynomial actions and their norm bound
are supplied. In the square-root theorem the finite identities use
induction on integer coefficients or exact rational matrices. Preservation
of a closed estimate follows by allocating the precision losses above
at each finite test.

\subsection{Why density matrices and the trace pairing suffice}

A density matrix is \(\rho\ge0\) with \(\operatorname{tr}\rho=1\).
If \(S=\sqrt\rho\), then \(\rho=SS^*\). Consequently
\[
 \operatorname{tr}(\rho E)=\sum_{j=1}^d (Se_j)^* E(Se_j),\qquad
 \sum_{j=1}^d\|Se_j\|_2^2=1 .
\]
Thus the Born expectation is real and in \([0,1]\) for an effect, and
is additive on sums of effects which are again effects. These identities
follow by interchanging finite sums, which also proves the cyclic trace
identity for compatible rectangular products.

\begin{theorem}[Representation from experimental additivity]\label{thm:q-effect}
Suppose \(f\) assigns a named number in \([0,1]\) to each effect, respects
equality, has \(f(I)=1\), and satisfies
\(f(E+F)=f(E)+f(F)\) whenever \(E+F\le I\). Then a unique density matrix
\(\rho\) represents it by \(f(E)=\operatorname{tr}(\rho E)\), in every
dimension including two. The result applies schematically to a supplied
evaluator with these laws.
\end{theorem}
\begin{proof}
First \(f(0)=0\). For \(A\ge0\) choose an integer \(n\ge1\) with
\(A\le nI\), and put \(g(A)=n f(A/n)\). Two choices \(n,m\) give the
same answer by adding \(A/(nm)\) successively \(m\) or \(n\) times.
Using one common bound for \(A,B,A+B\) now proves additivity of
\(g\) on the positive cone. Extend it to Hermitian \(H\) by
\(\ell(H)=g(H+kI)-k\), where \(k\) is an integer norm bound.
Independence of \(k\) follows from additivity and \(g(I)=1\).
One can also use \(k=0\) for positive \(H\). Taking separate
shifts for \(H,K\) and their sum as a shift for \(H+K\) gives
additivity of \(\ell\); it is positive and order preserving.
It follows that \(\ell(qI)=q\) for rational \(q\), including negative
rationals. If \(-qI\le H\le qI\) for positive rational \(q\), then
\(|\ell(H)|\le q\). Thus \(\ell\) is continuous and additive,
hence real-linear by rational approximation of the scalar. Quantitatively,
replacing a scalar by a rational within \(\eta\) changes the value on
a matrix of norm at most \(b\) by at most \(b\eta\).

Extend complex-linearly by
\(L(Z)=\ell((Z+Z^*)/2)+i\ell((Z-Z^*)/(2i))\).
This gives \(L(Z^*)=\overline{L(Z)}\). For matrix units \(E_{ij}\), set
\(\rho_{ji}=L(E_{ij})\). Expansion in these units gives
\(L(Z)=\operatorname{tr}(\rho Z)\), and determines every entry uniquely.
The adjoint identity makes \(\rho\) Hermitian,
\(x^*\rho x=L(xx^*)\ge0\), and \(\operatorname{tr}\rho=L(I)=1\).
All the applications of the original evaluator needed to construct
these entries are on explicit shifts and rescalings of the Hermitian
parts of finitely many matrix units. The additivity assumption concerns
all effects. Restricting it to orthogonal sharp projections would change
the theorem.
\end{proof}

For systems \(A,B\) of dimensions \(m,n\), use the lexicographic product
basis indexed by \((i,a)\) and define
\((S\otimes T)_{ia,jb}=S_{ij}T_{ab}\).
The partial trace has entries
\((\operatorname{tr}_B R)_{ij}=\sum_a R_{ia,ja}\).
It preserves trace and positivity: the quadratic form on \(x\) is the
sum of the quadratic forms of \(R\) on \(x\otimes e_a\).
It is characterized by
\(\operatorname{tr}(R(T\otimes I))=
\operatorname{tr}((\operatorname{tr}_B R)T)\), by testing matrix units.
Every density matrix has a purification on \(A\otimes A\):
the vector with entries \(v_{ia}=(\sqrt\rho)_{ia}\) has norm one and
\(\operatorname{tr}_2(vv^*)=\rho\).
Both statements follow by summing columns of \(\sqrt\rho\).
In particular, neither the extraction of eigenvectors nor the
normalization of possibly zero columns occurs in this construction.

\subsection{Channels, measurements and composite systems}

For a linear matrix map \(\Phi:M_d\to M_m\), the entries
\(\Phi(E_{ij})_{ab}\) specify its entire action. On block matrices
\(\operatorname{id}_k\otimes\Phi\) acts by applying \(\Phi\) to each
block. Complete positivity means positivity on every such positive
block matrix. A channel is a completely positive, trace-preserving map.

\begin{theorem}[Finite Choi and Kraus representations]\label{thm:q-choi}
The following descriptions of \(\Phi\) agree constructively:
complete positivity; positivity of the single \(dm\)-square matrix
\(J_{ia,jb}=\Phi(E_{ij})_{ab}\); and a representation
\[
 \Phi(X)=\sum_{s=1}^{dm}K_s X K_s^*,\qquad K_s:\mathbb C^d\to\mathbb C^m.
\]
Trace preservation is equivalent to \(\sum_s K_s^*K_s=I_d\).
\end{theorem}
\begin{proof}
Let \(w=\sum_i e_i\otimes e_i\). The positive matrix \(ww^*\)
has blocks \(E_{ij}\), so \((\operatorname{id}_d\otimes\Phi)(ww^*)=J\).
If \(J\ge0\), form \(R=\sqrt J\) by Theorem~\ref{thm:q-sqrt}, and set
\((K_s)_{ai}=R_{ia,s}\). Then
\[
 \bigg(\sum_s K_s E_{ij}K_s^*\bigg)_{ab}
 =\sum_s R_{ia,s}\overline{R_{jb,s}}=J_{ia,jb}.
\]
Linearity gives the formula on all matrices. Conversely on a positive
block matrix \(Y\), each conjugation
\((I_k\otimes K_s)Y(I_k\otimes K_s)^*\) is positive by a
quadratic-form test. Their sum is positive. Cyclicity gives
\(\operatorname{tr}\Phi(X)=\operatorname{tr}(X\sum_s K_s^*K_s)\).
Equality with \(\operatorname{tr}X\) for every matrix unit is exactly
the displayed condition. The construction uses all \(dm\) columns,
including zero columns, so its size is independent of any rank decision.
\end{proof}

A channel admits an isometric dilation: set
\(Vx=\sum_s K_sx\otimes e_s\). The normalization condition gives
\(V^*V=I\), and \(\operatorname{tr}_2(VXV^*)=\Phi(X)\) by the partial
trace formula. Successive channels with lists \(K_s,L_t\) have list
\(L_tK_s\), and independent channels have the list \(K_s\otimes L_t\).
These identities follow by distributing the sums. Their normalizations
also follow directly, e.g.\ \(\sum_{s,t}K_s^*L_t^*L_tK_s=I\).
Adding a prepared environment and subsequently discarding it therefore
has the same probability rule as the corresponding channel.

An instrument consists of lists \(K_{a,s}\) with
\(\sum_{a,s}K_{a,s}^*K_{a,s}=I\), grouped by a recorded outcome \(a\).
Its unnormalized output
\(\sigma_a=\sum_s K_{a,s}\rho K_{a,s}^*\) carries weight
\(p_a=\operatorname{tr}\sigma_a
=\operatorname{tr}(\rho E_a)\),
where \(E_a=\sum_s K_{a,s}^*K_{a,s}\). This is a finite probability
distribution since \(E_a\ge0\) and \(\sum_a E_a=I\).
Conversely, any such POVM (positive operator-valued measure) is realized
by \(K_a=\sqrt{E_a}\). Its dilation displays the POVM as a sharp
measurement on a larger system: \(E_a=V^*(I\otimes e_a e_a^*)V\).
For a projective measurement \(E_a=P_a\), \(P_a=P_a^*=P_a^2\),
with \(P_bP_a=0\) for distinct outcomes, this instrument repeats its
answer with certainty because the second branch is
\(P_bP_a\rho P_aP_b\). General instruments for a given POVM can have
different disturbances.

If a lower witness \(p_a\ge 2^{-h}\) is supplied, the state after
observing \(a\) is \(\sigma_a/p_a\). At denominator error
\(\eta\le2^{-h-1}\), a reciprocal computed from the approximate
denominator has error at most \(2^{2h+1}\eta\); this follows by taking
the difference of fractions. Even when no such witness is available,
the joint probability of seeing \(a\) and subsequently measuring an
effect \(F\) is unconditionally defined as
\(\operatorname{tr}(F\sigma_a)\).
For disjoint subsystems, an unrecorded operation on \(B\) leaves the
marginal on \(A\) invariant. Indeed cyclicity in the expectation of
\(T\otimes I\) replaces it by
\(\sum_s (I\otimes K_s)^*(T\otimes I)(I\otimes K_s)
=T\otimes I\). The matrix-unit characterization of partial trace gives
equality of the marginals.

\subsection{Discrimination, uncertainty and entanglement}

\begin{theorem}[Optimal binary discrimination with an effective tolerance]
\label{thm:q-helstrom}
Suppose the preparation is \(\rho\) with probability \(u\in[0,1]\)
and \(\tau\) with probability \(1-u\). Put
\(D=u\rho-(1-u)\tau\) and \(|D|=\sqrt{D^2}\).
The least upper bound on the success probabilities of binary measurements is
\[
             P_* = \tfrac12(1+\operatorname{tr}|D|).
\]
For positive rational \(\epsilon\), an explicit effect obtains at least
\(P_*-d\epsilon/2\).
\end{theorem}
\begin{proof}
Write \(B=|D|\), \(R=B+\epsilon I\). To invert a positive matrix with
\(bI\le R\le MI\), \(0<b\le M\) rational, form
\(M^{-1}\sum_{j=0}^{n}(I-R/M)^j\).
The factor is a positive contraction with bound \(q=1-b/M\);
apply Lemma~\ref{lem:q-forms} after scaling if \(q>0\), with the case
\(q=0\) immediate. The geometric tail is at most \(q^{n+1}/b\).
For \(0\le a\le1\), the inequality
\((1-a)^j\le1/(1+ja)\) follows by induction after multiplying the
denominators. Thus the tail is at most \(M/(b^2(n+1))\).
Multiplication of the partial sums by \(R\) telescopes to
\(I-(I-R/M)^{n+1}\), so the limit is a positive inverse. Here take
\(b=\epsilon\) and an integer upper bound for \(M\).

The commuting matrices \(D,B,R^{-1}\) give a self-adjoint contraction
\(T=D R^{-1}\): indeed \(R^2-D^2=2\epsilon B+\epsilon^2 I\ge0\),
so conjugation by \(R^{-1}\) gives \(T^2\le I\).
Hence \(I\pm T\ge0\) by Cauchy--Schwarz, and the commuting product
\(R(I\pm T)=B+\epsilon I\pm D\) is positive. Letting \(\epsilon\)
decrease to zero shows \(D_\pm=(B\pm D)/2\ge0\). Given any
effect \(E\), guess \(\rho\) on this effect and \(\tau\) otherwise.
The success probability is \(1-u+\operatorname{tr}(ED)\).
For positive \(X,Y\), \(\operatorname{tr}(XY)\) is nonnegative
because it equals \(\operatorname{tr}(\sqrt X\,Y\sqrt X)\).
Applying this to \((I-E)D_+\) and \(ED_-\) shows
\(\operatorname{tr}(ED)\le\operatorname{tr}D_+\). As
\(\operatorname{tr}D=2u-1\), this yields \(P_*\).

Use \(E_\epsilon=(I+T)/2\) for the reverse estimate. Its deficit is
\(\tfrac12\operatorname{tr}(B-B^2R^{-1})
=\tfrac{\epsilon}{2}\operatorname{tr}(B R^{-1})\).
The product \(BR^{-1}\) is positive, and
\(I-BR^{-1}=\epsilon R^{-1}\ge0\), so its trace is at most \(d\).
For target deficit \(2^{-p}\) one may take \(\epsilon=2^{-p}/d\).
This realizes the stated least-upper-bound property with a uniform
primitive-recursive approximation to the measurement.
\end{proof}

For equal priors the distinguishability
\(\tfrac12\operatorname{tr}|\rho-\tau|\) cannot increase under a channel.
Every output effect \(E\) pulls back to the input effect
\(\sum_s K_s^*E K_s\) with the same success probability, by cyclicity
and the channel normalization. Taking the explicit near-optimal output
tests of the theorem proves the inequality to arbitrary precision.

For a self-adjoint observable \(A\), put
\(\langle A\rangle=\operatorname{tr}(\rho A)\),
\(A_0=A-\langle A\rangle I\), and
\(\operatorname{Var}_\rho A=\langle A_0^2\rangle\).
The Hermitian form \(h(S,T)=\operatorname{tr}(\rho S^*T)\) is positive
because \(h(S,S)=\sum_j\|S\sqrt\rho\,e_j\|_2^2\).
Lemma~\ref{lem:q-forms} applied to \(A_0,B_0\) gives the full
Schr\"odinger--Robertson inequality
\[
 \operatorname{Var}_\rho A\,\operatorname{Var}_\rho B\ \ge \left(\tfrac12\langle A_0B_0+B_0A_0\rangle\right)^2+
 \left(\tfrac1{2i}\langle AB-BA\rangle\right)^2 .
\]
Indeed the displayed terms are the squares of the real and imaginary
parts of \(h(A_0,B_0)\). There is no division by a variance.

\begin{example}[A Bell experiment with an exact separation]\label{ex:q-bell}
For the two real Pauli matrices
\[
 Z=\begin{pmatrix}1&0\\0&-1\end{pmatrix},\qquad
 X=\begin{pmatrix}0&1\\1&0\end{pmatrix},
 \qquad v=(e_1\otimes e_1+e_2\otimes e_2)/\sqrt2 ,
\]
choose \(A_0=Z,A_1=X\), \(B_0=(Z+X)/\sqrt2\),
\(B_1=(Z-X)/\sqrt2\). Each observable is a self-adjoint involution
since \(ZX+XZ=0\); its two effects are \((I\pm A)/2\).
The CHSH observable is
\[
 C=A_0\otimes(B_0+B_1)+A_1\otimes(B_0-B_1)
   =\sqrt2(Z\otimes Z+X\otimes X).
\]
Both summands inside the parentheses fix \(v\), so \(v^*Cv=2\sqrt2\).
In a product state the value is
\(a_0(b_0+b_1)+a_1(b_0-b_1)\), where each expectation lies in
\([-1,1]\). Its absolute value is at most
\(|b_0+b_1|+|b_0-b_1|=2\max(|b_0|,|b_1|)\le2\);
the middle real identity follows by rational approximation from the
four rational sign cases. The bound \(2\) passes to all finite
convex mixtures and their effective limits. Thus the Bell state is
entangled, with a positive margin. For arbitrary self-adjoint
involutions on the two sides, expansion gives
\(C^2=4I-[A_0,A_1]\otimes[B_0,B_1]\), with
\([S,T]=ST-TS\). Each commutator is a difference of two unitaries.
Its tensor product here expands as four signed unitaries; hence
\(\|C^2\|\le8\), giving \(\|Cx\|_2^2=x^*C^2x\le8\|x\|_2^2\).
Thus the stated experiment attains the quantum bound for these
measurements. Its coefficients require only rational arithmetic and
a square root of the fixed positive rational \(2\).
\end{example}

\paragraph{Finite adaptive protocols.}
Suppose each internal node $v$ of a finite tree specifies an instrument with matrices $K_{a,s}^{(v)}$ and successor $va$. The choice of subsequent instrument may depend on all previous recorded outcomes. Give each leaf its final acceptance effect $F_v$, with $0$ and $I$ corresponding to rejection and acceptance. Working back toward the root, compute
\[
 F_v=\sum_{a,s}(K_{a,s}^{(v)})^* F_{va} K_{a,s}^{(v)}.
\]
Quadratic-form testing and normalization give $0\le F_v\le I$. Repeated finite cyclic traces show that $\operatorname{tr}(\rho F_{\rm root})$ equals the sum of the joint Born weights of all accepted histories (or the final weighted acceptances). This evaluates an adaptive experiment even if some histories have zero probability. If all child effects are determined within operator error $\eta$, the parent incurs at most $\eta$ in inherited error: compress the diagonal operator with blocks $F_{va}$ along the isometry with components $K_{a,s}^{(v)}$. For tree depth $D$, choose local matrix arithmetic error at most $\delta_{k+3}/(1+D)$ at every node. The backward evaluation then contributes at most $\delta_{k+3}$, independently of how many histories have small weight. Finite sums at a node use the earlier dimension- and fan-out-dependent guard digits. Reserve a separate precision request for the initial trace. For variable tree size the evaluation is PR on a finite tree code with postorder numbering and a joint sampler for its node data.

\section{A countably infinite Hilbert space with quantitative names}
\label{sec:infinitehilbert}

Fix the discrete basis $e_0,e_1,\ldots$. A finite vector is a finite list of
Gaussian rationals, with trailing zeroes ignored. Write
$\delta_k=2^{-k}$, $P_N$ for restriction to $0\le i<N$, and
$\lambda(C)=\min\{h\le C:2^h\ge C\}$ for positive integral $C$.
Ceil a positive rational before applying $\lambda$. Squared norms of finite
vectors are rational. Thus even a test $\|u-v\|\le r$ for nonnegative
rational $r$ is a rational test. Inner products are conjugate-linear in the
first slot. All algorithms below apply to fixed PR presentations or uniformly
to families carrying PR algorithms, in the schematic sense of iPRA.

\subsection{Vectors, tails and completeness}

A name consists of an integer $M\ge1$ and a PR function $s\mapsto x_s$ of
finite vectors with
\[
 \|x_s\|\le M,\qquad \|x_s-x_t\|\le\delta_s\quad(t\ge s).
\]
Equality of two names means $\|x_s-y_s\|\le 2\delta_s$ for every $s$.
This relation is transitive: compare both names at an arbitrarily late
stage with an intermediate name, then return to stage $s$. The resulting
bound is $2\delta_s+4\delta_t$ for every $t\ge s$; a strict rational
violation of $2\delta_s$ supplies a finite $t$ excluding it. Denote the
setoid by $\ell^2_{\rm PR}$. An estimate such as $\|x-x_s\|\le\delta_s$
means the corresponding estimates at finite stages with their rounding
allowances. In computations one can target error $\delta_{k+2}$ at stage $k$:
two such approximations to the same vector differ by at most $\delta_k$
at stages $k$ and $t\ge k$.

The finite Cauchy--Schwarz inequality follows, including the zero-vector
case, from the polynomial identity
\[
 \|a\|^2\|b\|^2-|\langle a,b\rangle|^2
   =\sum_{i<j}|a_i b_j-a_j b_i|^2.
\]
It gives the triangle inequality by expanding a square. Addition,
multiplication by a named scalar, the norm, and the inner product pass to
names. For example, when inputs have bounds $M,N$, the inner-product
error made by using both stage-$s$ vectors is at most $(M+N)\delta_s$;
the norm error is at most $\delta_s$ before rational square-root rounding.
For scalar products use
$\|a x-b y\|\le |a|\|x-y\|+|a-b|\|y\|$. A stage
$k+3+\lambda(1+M+N)$ therefore suffices for the inner product and
for the analogous scalar-vector product with these envelopes. Finite norm
rounding uses bisection of a nonnegative rational square root on $[0,M]$.
The finite identities extend by these bounds, including positivity,
conjugate symmetry and the parallelogram law. If $\|x\|=0$ then
$\|x_s\|\le\delta_s$ and $x=0$.

\begin{proposition}[Two presentations of the basis expansion]
A vector name gives uniformly named coordinates and a PR cutoff $L(k)$
such that every finite section of its tail beyond $L(k)$ has norm at
most $\delta_k$. Conversely, uniformly named coordinates with this
finite-tail certificate and a fixed bound $M$ on the norms of all finite
initial sections give a vector name.
\end{proposition}
\begin{proof}
Read a coordinate from $x_s$ with error at most $\delta_s$. Take $L(k)$
equal to the length of $x_k$. Every finite selection of entries beyond it
is a contraction of $x-x_k$. Conversely take $N=L(k+3)$ and approximate the
first $N$ coordinates to complex absolute error
$\delta_{k+3+\lambda(1+N)}$ each. The resulting vector $q_k$ has error at
most $\delta_{k+3}$ in that section, since the norm is at most the sum
of coordinate errors. For $s\ge k$, compare $q_k,q_s$ inside the longer
of their two sections; each incurs its rounding error and at most its
certified tail error. Their difference has norm at most $4\delta_{k+3}$.
The integer $M+1$ bounds their norms. Coordinate recovery follows by
comparing in a section containing the desired index, using the tail estimate
even when that index lies beyond the chosen cutoff. Reconstruction of a
given name is equal to the original by the same estimates. A monotone
cutoff may always be obtained by taking finite running maxima.
\end{proof}

The same estimates give the relevant completeness principle. Suppose
$y^{(j)}$ is a uniformly named family, and a PR function $m$ certifies
$\|y^{(a)}-y^{(b)}\|\le\delta_r$ for $a,b\ge m(r)$.
Replace $m$ by running maxima. For stage $k$ take stage $k+3$ of
$y^{(m(k+3))}$. Two such samples differ by at most $3\delta_{k+3}$.
Their norms are bounded, for example, by the bound of $y^{(m(0))}$
plus $2$. Comparison against a later sample proves convergence with the
given modulus. This diagonal construction uses a supplied joint PR sampler.
A collection of individually PR names need not have one.

The tail certificate carries mathematical information. Let $h(e,j)$ be the
decidable predicate that machine $e$ first halts at step $j$. The sequence
$a_j=h(e,j)$ has at most one nonzero coordinate, and every finite partial
sum of squared moduli is at most one. A uniform computable cutoff $L_e(1)$
would decide halting by simulating up to that cutoff: a later $1$ would give
a finite tail of norm one. This external reduction also excludes a PR rule
producing such cutoffs from the coordinate algorithm and norm upper bound.
It explains precisely why those two pieces of data alone do not give the
above Hilbert-space presentation.

\subsection{Bounded maps and tensors}

A bounded-map presentation for $A$ gives uniformly named columns
$c_j=Ae_j$ and an integer $B\ge1$ with the certificate
\[
 \Big\|\sum_j q_jc_j\Big\|\le B\|q\|
 \qquad\text{for every finite Gaussian-rational vector }q .
 \tag{H1}\label{eq:H1}
\]
The inequality expands to universal precision tests. For an actual
adjointable map one supplies a second bounded-map presentation with columns
$d_i$, and the identities $(c_j)_i=\overline{(d_i)_j}$.
When needed enlarge $B$ to cover both bounds. There is no operation
of searching for a norm or an adjoint in these definitions.

\begin{theorem}[Quantitative extension and adjoint]
These data determine an extensional linear map on vector names, bounded by
$B$. The second presentation acts as its Hilbert-space adjoint.
\end{theorem}
\begin{proof}
For target stage $k$ take $s=k+3+\lambda(B)$, $q=x_s$, and
$C=1+\lceil\sum_j(|\Re q_j|+|\Im q_j|)\rceil$. Return the finite sum
of $q_j$ times the stage-$k+3+\lambda(C)$ columns. Rounding the columns
costs at most $\delta_{k+3}$. Changing from $x_s$ to any subsequent
input stage costs at most $B\delta_s\le\delta_{k+3}$ by
\eqref{eq:H1}. Thus the stage vectors are regular, have bound $BM+1$,
and represent the unique bounded extension, since finite input approximants
are dense with a known rate. Addition and scalar multiplication follow by
comparison on finite inputs and taking the prescribed error to zero.
Pairing the columns on two finite inputs gives
$\langle y,Ax\rangle=\langle A^*y,x\rangle$.
The inner-product error bound extends this identity to names.
In particular, if all columns vanish, the whole map vanishes.
\end{proof}

Columns of $AD$ are $A(De_j)$, the norm bound is the product of the bounds,
and $(AD)^*=D^*A^*$. For a sum take a sum of bounds. An isometry can be
specified by the finite-input identity $\|Aq\|^2=\|q\|^2$; a supplied
unitary additionally has a two-sided inverse, for example by certifying
$A^*A=AA^*=I$ on the columns. These presentations are closed under all
three constructions just given, with nested samplers primitive recursive.

Code a pair of indices using a fixed PR bijection with PR inverses, e.g.\
diagonal enumeration of $\mathbb N^2$ with bounded inversion. The tensor
space consists of regular names of finitely supported arrays in the pair
index. Tensoring finite vectors gives the identities
\[
 \|x\otimes y\|=\|x\|\|y\|,\qquad
 \langle x\otimes y,u\otimes v\rangle
   =\langle x,u\rangle\langle y,v\rangle .
\]
They extend to names because approximating the two factors incurs error
at most $(M+N)\delta_s$ for stage bounds $M,N$. If $A$ is bounded, then
$A\otimes I$ acts on each second-index column of a finite array. Sum the
squares of \eqref{eq:H1} over those columns to obtain bound $B$ on the
array norm. Its columns are $Ae_i\otimes e_j$; the extension theorem
therefore applies. The argument for $I\otimes D$ interchanges the indices.
Composition gives $A\otimes D$ with the product bound; the adjoint and
permutation identities hold on basis arrays and extend. This realizes the
completed separable tensor product with the given bases.

\subsection{Nuclear operators and an intrinsic trace}

For vectors $u,v$ write $[u,v]x=u\langle v,x\rangle$. A nuclear
presentation consists of uniformly named $u_j,v_j$, nonnegative PR
rationals $m_j\ge\|u_j\|\|v_j\|$, an integer $C$ bounding all finite
partial sums of the $m_j$, and a PR tail function $T$ with
\[
  \sum_{a\le j<b}m_j\le\delta_r\qquad(b\ge a\ge T(r)).
 \tag{H2}\label{eq:H2}
\]
It presents $Sx=\sum_j[u_j,v_j]x$. Two presentations represent the same
operator when all matrix entries agree. An upper bound for the classical
trace norm is supplied by $C$; no minimization over decompositions is used.
Finite-rank rational operators, and PR nuclear series equipped with the
displayed certificates, are included.

\begin{theorem}[Trace and bounded observables]
$S$ and $S^*$ have bounded presentations with bounds $\max(1,C)$.
For any bounded-map presentation $A$ with bound $B$, the scalar
\[
  \mathcal T_A(S)=\sum_j \langle v_j,A u_j\rangle
\]
is a regular complex name depending only on the operator represented by
$S$. Thus $\operatorname{Tr}S=\mathcal T_I(S)$ is intrinsic, and
$|\mathcal T_A(S)|\le BC$. If $A$ has a presented adjoint, then
$\operatorname{Tr}(AS)=\mathcal T_A(S)=\operatorname{Tr}(SA)$.
\end{theorem}
\begin{proof}
For $\|x\|\le X$, each term of $Sx$ has norm at most $Xm_j$.
Equation \eqref{eq:H2} provides the Cauchy modulus $T(r+\lambda(X))$
when $X\ge1$. Apply the diagonal construction to finite partial sums;
swap $u_j,v_j$ for $S^*$. Bounds and the adjoint identity follow first
for partial sums and then at every precision.

The scalar tail is at most $B\delta_r$. For an explicit stage-$k$
approximation choose $r=k+3+\lambda(B)$, $J=T(r)$.
Let $U_j,V_j\ge1$ be integer envelopes for $u_j,v_j$.
Take
\[
 h=k+3+\lambda\!\left(1+\sum_{j<J}(BU_j+V_j+1)\right).
\]
Pair finite stage-$h$ samples for $Au_j,v_j$. The first sample has
distance at most $\delta_h$ from a vector of norm at most $BU_j$, hence
the summed inner-product error is at most
$\sum_{j<J}(BU_j+V_j+1)\delta_h$. Regularity follows.

To see independence of the expansion, use a finite basis trace
$D_N=\sum_{i<N}\langle e_i,A S e_i\rangle$. Distributing the finite sum
over finite partial nuclear sums and using their tail bounds gives
$D_N=\sum_j\langle P_Nv_j,A u_j\rangle$ for each fixed $N$.
For $j\ge J$ the total absolute difference from $\mathcal T_A(S)$
is bounded by $B\delta_r$. For $j<J$ take
$h'=k+3+\lambda(1+B\sum_{j<J}U_j)$ and choose $N$ large enough
to contain the supports of the first $J$ samples $(v_j)_{h'}$.
Then $\|(I-P_N)v_j\|\le\delta_{h'}$, and the difference over the
head is at most $B\sum_{j<J}U_j\delta_{h'}$. Both estimates persist
for larger $N$. Two presentations with the same matrix entries have the
same columns, since a named vector with every coordinate zero has norm
zero by its effective tails. Thus they have identical $D_N$; take the
larger of their two cutoffs to prove equality of their traces. Finally
$AS$ has expansion $[Au_j,v_j]$ and $SA$ has expansion
$[u_j,A^*v_j]$. Their majorants are $Bm_j$, with the tail request shifted
by $\lambda(B)$, and the adjoint identity equates their traces.
\end{proof}

\subsection{Density states, discarding a system and observing an outcome}

A convenient positive-state presentation is a unit vector $w$ on a
system and a countable ancillary basis. Finite ancillary bases are included
by padding with zeroes. Write $w_j$ for the system vector obtained by
keeping ancilla index $j$. These slices have joint names by contraction.
Define
\[
  \rho_w=\sum_j[w_j,w_j],\qquad
  \omega_w(A)=\langle w,(A\otimes I)w\rangle .
\]
The presentation specifies the probabilities of experiments on the system.
Distinct amplitudes represent the same mixed state if their density
operators have identical matrix entries.

\begin{proposition}[Density and reduction]
$\rho_w$ has a nuclear presentation, is positive, and has trace one.
For every presented bounded $A$, $\omega_w(A)=\mathcal T_A(\rho_w)$.
Discarding a factor of a multipartite system amounts to moving its
index into the ancilla. This gives a presentation-independent operation
on these states.
\end{proposition}
\begin{proof}
For every finite set of indices, the sum of slice norms squared is at
most $\|w\|^2$ by orthogonal coordinate restriction. Choose a rational
$m_j$ in $[\|w_j\|^2,\|w_j\|^2+\delta_{j+2}]$ by approximating
the square with two extra bits and adding its error. The budget is at
most $1+\sum_j\delta_{j+2}\le 2$. For target $r$ let $J$ exceed
every ancilla index occurring in the finite sample $w_{r+2}$, and
increase it to at least $r+2$. Every finite slice tail beyond $J$
has total squared norm at most $\delta_{r+2}^2$, by contraction
applied to $w-w_{r+2}$. The rounding slacks sum to at most
$2^{-J-1}$. Their sum with $\delta_{r+2}^2$ is below $\delta_r$.
Thus \eqref{eq:H2} holds. Also the finite ancilla projection up to $J$
approximates $w$ in norm to within $\delta_{r+2}$, so summing all
squared slice norms gives exactly one.
For finite projections the pairing formula is just summing over
orthogonal slices. The norm bound for $A\otimes I$ passes it to the
limit. In particular the trace is one and
$\langle x,\rho_w x\rangle=\sum_j|\langle w_j,x\rangle|^2\ge0$.
If an additional system factor is discarded, rearrange the two ancillary
indices with the pairing bijection. For every remaining bounded test $A$
the new expectation equals the old expectation of $A\otimes I$ by
the pairing formula. Testing matrix units recovers every density entry,
which proves independence of the initial purification.
\end{proof}

\begin{theorem}[Every presented density has a purification]
Suppose $S$ has a nuclear presentation, is self-adjoint and positive,
and has trace one. One can construct a uniform unit-vector name $w$ on
system $\otimes$ ancilla such that $\rho_w=S$.
\end{theorem}
\begin{proof}
Take $L=1+C$, so $0\le S\le L^2I$.
The bounded-map polynomial algorithm of Theorem~\ref{thm:q-sqrt}
applied to columns constructs a positive self-adjoint $R$ with $R^2=S$.
For a column target of $\delta_h$ one may take
$m=L^22^{2h+4}$ terms and use the bounded-map extension for the
finite polynomial to a further $\delta_{h+2}$ of error. In detail, $X=I-S/L^2$ is contractive. Successive applications to a
finite vector can each be rounded to $\delta_{h+3}/[L(m+1)]$; the
$j$th iterate then has error at most $j$ times this budget. The absolute
nonconstant coefficients of the square-root polynomial sum to at most
$L$. Thus the variable number of powers uses one PR iteration on finite
vectors and its total error is at most $\delta_{h+3}$.
Set the desired $j$th ancillary slice to $R e_j$. Then
\[
 \sum_{j<N}\|R e_j\|^2
    =\sum_{j<N}\langle e_j,S e_j\rangle=D_N .
\]
The basis-trace proof with $A=I$ gives a PR cutoff $J(r)$ past which
$|D_N-1|\le\delta_r$. The left side is nondecreasing, and its finite
tails from any $N\ge J(r)$ are bounded by $\delta_r$, since each
partial sum is at most its named limit one. At stage $k$ use
$N=J(2k+6)$; assemble approximations to the first $N$ column vectors
of $R$ to error $\delta_{k+3+\lambda(1+N)}$ each.
The omitted slice vector has norm at most $\delta_{k+3}$ in every
finite section, and the finite rounding error is at most $\delta_{k+3}$.
These estimates prove regularity of the assembled name and its norm
equals one by the displayed identity. For a matrix entry of $\rho_w$
the finite slice sum is the corresponding entry of $R P_N R$ by
self-adjointness. On the input $e_i$ its difference from $R^2$ has
norm at most $L\|(I-P_N)R e_i\|$, with PR cutoff from the named
column $R e_i$. Passing to the limit gives $\rho_w=R^2=S$.
\end{proof}

A mixed preparation with PR unit vectors $x_j$ and nonnegative named
probabilities $p_j$ summing to one with PR tail modulus $T$ is included:
use slices $w_j=\sqrt{p_j}\,x_j$. A square root of $p\in[0,1]$ is
obtained by approximating $p$ within $\delta_{2k+6}$, clipping to
$[0,1]$, and bisecting the rational approximation's square root to
$\delta_{k+3}$; $|\sqrt a-\sqrt b|\le\sqrt{|a-b|}$ controls the
first error. Choose $J=T(2k+6)$ and approximate the first $J$ slices
to $\delta_{k+3+\lambda(1+J)}$ each. The assembled error is at most
$\delta_{k+3}$ from the finite approximations and $\delta_{k+3}$ from
the omitted slices. Values $p_j=0$ are allowed. This construction does
not require an eigenbasis for a density operator.

An operation is supplied by a presented isometry $V$ from the input to
output $\otimes$ environment, together with its adjoint when a Heisenberg
map is desired. On a purification apply $V\otimes I$ and move the new
environment index into the ancilla. Its trace is one by isometry. The
output expectation of $A$ is the input expectation of
$V^*(A\otimes I)V$. Consequently the state operation is extensional
and affine, since matrix entries are tested by matrix units. At every
finite matrix level a positive presented test remains positive under
this Heisenberg transformation: for an extra finite factor take the
quadratic form on $(I\otimes V)z$, after permuting the environment
factor. Positivity of a test tensored with the environment identity
follows by summing the nonnegative quadratic forms on finite slices
and applying the explicit tail estimate. The identity test is fixed by
$V^*V=I$. This proves complete positivity and normalization for these
dilations. With only a forward map the pure-vector sampling construction
still works; the adjoint presentation supplies the stated observable
transformation and the argument for extensionality on mixed preparations.

For a countable measurement use a distinguished outcome index in $V$'s
environment. Let $z=(V\otimes I)w$, permuting indices so that the
outcome comes last, and $z_a$ its outcome-$a$ slice. The Born rule and
conditional state become
\[
   p_a=\|z_a\|^2,\qquad
   \rho_{z_a/\sqrt{p_a}}\quad\text{given a witness }p_a\ge\delta_b>0 .
\]
Here $z_a$ retains the output and any unobserved ancillary factors.
The unnormalized slice always exists. For the quotient, $\sqrt{p_a}$
has certified lower bound $\delta_b$, and reciprocation on this region
has Lipschitz bound $2^{2b}$. Rational estimates may first be clipped
to $[\delta_b,1]$. Every PR-decidable event has probability
$\|(I\otimes P_E)z\|^2$, where the projection simply filters the
finite sampled array by its outcome. Orthogonality gives finite
additivity and conditioning on a disjoint union gives the sum of its
unnormalized branch states, as is seen by the slice expansion. The
cutoff determined from the largest outcome in $z_{r+2}$ bounds the
remaining mass by $\delta_{r+2}^2$.
This proves countable normalization and supplies its rate even for
an infinite-outcome device. More generally a jointly PR disjoint family
of events admits effective summation to its union when there is a PR
bound locating the part containing each accepted outcome among the
finitely many retained indices; such a bound is automatic for the
singleton outcome decomposition used here.

\paragraph{Arithmetic content of the estimates.}
Support lengths, maxima of finite lists, bisection depths and all sampling
indices above are PR in their numeric parameters. Finite Pythagorean
identities, rearrangements and norm estimates can be verified as
integer or rational finite-array statements after clearing positive
denominators. Their loops use quantifier-free induction with the free
precision, cutoff and input parameters fixed. Uses of equalities of
completed objects compare arbitrary output precisions and then substitute
the displayed finer precisions. In particular the completion argument
and the trace proof do not perform induction over a quantified
convergence predicate.

\section{A spectral calculus built from finite measurements}\label{sec:spectrum}
For a bounded self-adjoint $A$ with a supplied integer bound $L\ge1$, write $E=(I+A/L)/2$. Then $0\le E\le I$. Conversely a point $u\in[0,1]$ represents the reading $L(2u-1)$. We first solve the measurement problem for $E$. An observable can then be tested with a continuous detector response $f$. No numerical spectrum, eigenbasis or spectral-gap decision enters this construction.

\begin{lemma}[Binomial measurements and the polynomial estimate]\label{lem:binomial}
For $N\ge1$, let $W_{Nj}={N\choose j}E^j(I-E)^{N-j}$ for $0\le j\le N$. They are effects summing to $I$. If a complex polynomial $p$ with rational coefficients satisfies $|p(t)|\le K$ on rational $t\in[0,1]$, then $\|p(E)\|\le K$. A real polynomial nonnegative on this interval gives a positive operator.
\end{lemma}
\begin{proof}
Products of commuting positive operators are positive by Theorem~\ref{thm:q-sqrt}; its proof by polynomials uses only Hilbert-space estimates. The binomial formula sums the products to $I$. Define $Vx=(\sqrt{W_{Nj}}x)_{j=0}^N$. Since $\sum W_{Nj}=I$, $V$ is an isometry, with adjoint the finite row $(\sqrt{W_{Nj}})_j$. For any complex list $c_j$, compression of its diagonal action yields
\[
 \left\|\sum_j c_jW_{Nj}\right\|\le\max_j |c_j|.
\]
In particular nonnegative coefficients give a positive operator. All maxima used in bounds may be replaced by given rational upper bounds.

For $p(t)=\sum_{a=0}^{d}p_at^a$, $N\ge\max(1,d)$, the exact expansion of $p(E)$ in $W_{Nj}$ has coefficients
\[
 d_{Nj}=\sum_a p_a\frac{(j)_a}{(N)_a},\qquad (j)_a=j(j-1)\cdots(j-a+1).
\]
Here $(j)_0=1$ and $(j)_a=0$ for $j<a$. Indeed ${N\choose j}(j)_a/(N)_a={N-a\choose j-a}$ when $j\ge a$, and summing after extracting $E^a$ proves the identity. For a population of $N$ labels with $j$ successes, these two quantities, $(j)_a/(N)_a$ and $(j/N)^a$, are probabilities of all successes in $a$ draws without and with replacement. Couple them by keeping the independent draws when they are distinct and using a fresh uniform distinct tuple otherwise. The conditional distribution on distinct tuples is uniform, including the replacement on the collision event; this can equivalently be checked by counting ordered tuples without any division by the probability of a collision. A union bound over the pairs gives probability of collision at most $a(a-1)/(2N)$. Thus
\[
  |d_{Nj}-p(j/N)|\le C_p/N,\qquad C_p=\tfrac12\sum_a |p_a|a(a-1).
\]
A rational upper bound using the absolute values of real and imaginary parts suffices. Hence $\|p(E)\|\le K+C_p/N$. Given a rational tolerance, choose $N$ by division by that tolerance. For real nonnegative $p$ the same estimate gives $p(E)\ge-(C_p/N)I$. Remove the slack in both estimates as in Section~\ref{sec:arithmetic}. This proof is a finite counting argument, including at end points and multiple spectral values.
\end{proof}

\begin{theorem}[Continuous functional calculus and the Born distribution]\label{thm:calculus}
Supply $f$ on rational points in $[0,1]$ by a joint PR value sampler, an integer $M\ge1$ bounding $|f|$, and a PR modulus $\mu$ such that $|u-v|\le2^{-\mu(r)}$ implies $|f(u)-f(v)|\le2^{-r}$. There is a unique complex-linear sup-norm contractive unital $*$-homomorphism from these functions (identified by their values on rationals) to presented operators taking $u\mapsto u$ to $E$. Its value $f(E)$ admits polynomial approximations with a PR operator-error bound. Each normalized state $\omega$ from the previous sections defines a positive normalized functional $f\mapsto\omega(f(E))$. The atomic distributions placing mass $\omega(W_{Nj})$ at $j/N$ approximate it at an explicit rate for each $f$.
\end{theorem}
\begin{proof}
Put $B_N f(t)=\sum_j f(j/N){N\choose j}t^j(1-t)^{N-j}$. For each rational $t$, these weights are a probability distribution with mean of $j/N$ equal to $t$ and variance $t(1-t)/N$. The identities follow from the falling-factorial calculation for $a=1,2$ (or directly when $N=1$). For any $\eta=2^{-\mu(r)}$, the weight outside $|j/N-t|\le\eta$ is at most $1/(4N\eta^2)$ by summing the squared deviations. It follows that
\[
 \sup_{t\in\mathbb Q\cap[0,1]}|B_N f(t)-f(t)|\ \le\ 2^{-r}+\frac{M}{2N\eta^2}.
\]
A supremum sign here abbreviates a common bound at every rational point. For target $k$ take $r=k+4$, $N\ge\max(1,\lceil M2^{r+2\mu(r)}\rceil)$, and replace each $f(j/N)$ by a rational complex approximation within $2^{-r}$. This adds at most $2^{-r}$ in uniform error and in compressed-operator error. Call the resulting rational polynomial $p_k$. Its uniform distance from $f$ is at most $3\delta_{k+4}$. By Lemma~\ref{lem:binomial}, $p_k(E)$ is uniformly Cauchy with an explicit rate. On each finite vector it gives a Hilbert name. To compute its action to target stage $k$, use polynomial $p_{k+2+\lceil\log_2(1+M_x)\rceil}$ for input norm bound $M_x$, then evaluate its action with arithmetic error at most $\delta_{k+3}$. Column extension and conjugate coefficients provide an adjoint. The same estimates give norm bound $M$.

Any other certified approximating polynomials give the same value by the polynomial estimate. In particular constant and identity functions have the indicated values. Sums, products and conjugates of approximating polynomials converge with the usual sum and product estimates, so evaluation preserves those operations. Two functions at uniform distance $b$ give outputs at distance at most $b$. Uniqueness follows by approximating at each requested tolerance by one $p_k$. If $f$ is real and nonnegative, its sampled weights may be clipped to nonnegative rationals, giving positive approximants. Evaluating a normalized state changes no operator-error bound: in a purification this is Cauchy--Schwarz on a unit vector. This proves both the probability assertion and its rate. Signed and complex detector responses are expectations, while $0\le f\le1$ gives an effect. A general bounded $A$ is handled by the indicated affine change of variable with a rescaled modulus.
\end{proof}

\subsection{Sharp questions and their extra data}
Different physical questions about the same $A$ can require different information. Suppose first that one is given a rational threshold $a$, a rational $L\ge1$ with $\|A-aI\|\le L$, and a rational $0<g\le1$ certifying $B^2\ge g^2I$ for $B=(A-aI)/L$. There is a sharp two-outcome spectral measurement separating the two signs. Indeed $X=I-B^2$ is positive with $\|X\|\le1-g^2=q$. Put
\[
 S=\sum_{n\ge0} 4^{-n}{2n\choose n}X^n,\qquad J=BS,\qquad P_{\pm}=(I\pm J)/2 .
\]
The tail after degree $m$ is at most $q^{m+1}/g^2\le 1/((m+1)g^4)$ by the geometric inequality used in Theorem~\ref{thm:q-helstrom}. For $g=1$ only the constant term is needed. The scalar coefficients $c_n$ obey $(n+1)c_{n+1}=(n+\tfrac12)c_n$. Comparing coefficients in $(1-z)F'=F/2$, $F(0)=1$, gives $(1-z)(F^2)'=F^2$ and all coefficients of $F^2$ are one. For each fixed degree these are finite convolution identities. The geometric bounds permit their use at $X$, giving $S^2(I-X)=I$. Thus $J^*=J$, $J^2=I$, and the $P_\pm$ are complementary commuting orthogonal projections. Since $BJ=B^2S$ is positive and commutes with the projections, $B$ is positive on $P_+$ and negative on $P_-$. More precisely $(BJ)^2=B^2$, so $BJ$ is the unique positive root of $B^2$. This sign decomposition is canonical and supplies the gap-certified spectral projections without selecting an eigenspace basis.

A single state's probability of an interval may exist with less information. For $0<c<1$ define clipped linear ramps $l_\eta,u_\eta$ for the event above $c$: $l_\eta$ rises from zero at $c$ to one at $c+\eta$, and $u_\eta$ rises from zero at $c-\eta$ to one at $c$. Their restrictions to $[0,1]$ are continuous effects. If a supplied PR schedule $h(k)$ proves $\omega((u_\eta-l_\eta)(E))\le\delta_k$ at $\eta=2^{-h(k)}$, the two expectations enclose a unique probability with this error. Use running maxima of $h$ to get nested intervals and sample their midpoints to two extra bits. Every continuous effect between these brackets has its expectation in the same interval. This constructs a sharp-event probability in that state. It does not purport to produce a projection on all states. Section~\ref{sec:particle} treats a different source of information: explicitly represented spatial intervals give projections directly.

\begin{proposition}[A uniform sharp cutoff is impossible]\label{prop:cut}
Even for one-dimensional uniformly PR self-adjoint matrices there is no computable general rule returning the probability of the strictly positive spectral half-line in the unit state.
\end{proposition}
\begin{proof}
Given machine $e$, let its real sampler at $k$ simulate it for $k$ steps. Return $2^{-j-1}$ if it first halted at step $j\le k$, and zero otherwise. This is jointly PR, bounded by one, and satisfies the regularity inequality: a later first halt changes the answer by at most $2^{-k-1}$. Take the resulting nonnegative scalar as the one-by-one matrix. In the ordinary semantics its positive spectral probability is exactly one when the machine halts and zero otherwise. Computing it to within $1/4$ would decide halting. This external impossibility proof applies even to an algorithm receiving the explicit PR source for the sampler. It precludes a total sound PR constructor in iPRA. It does not preclude the continuous response, gap or state-specific constructions above.
\end{proof}

\section{Unbounded observables and the Stone correspondence}\label{sec:stone}
An energy bound for a Hamiltonian usually fails, even on unit vectors. A useful input is the ability to solve $(H\pm iR)y=x$ at controlled error. The inverses are everywhere bounded, while their ranges remember the domain. Constructive Stone theorems have precedents in Bishop-style analysis \cite{ye}. Here both directions are equipped with first-order PR sampling and error bounds.

\subsection{Input and bounded exponential}
A presented graph is a linear relation on named pairs $(x,y)$ specified by arithmetical tests on their samples, functional up to Hilbert equality; write $y=Hx$. A graph member supplies both names and witnesses of the graph laws. We require symmetry $\langle Hx,z\rangle=\langle x,Hz\rangle$ on its domain. For each integer $R\ge1$ require jointly PR bounded maps $R_+(R),R_-(R)$ with
\[
 (H\pm iR)R_\pm(R)x=x,\qquad R_\pm(R)(H\pm iR)y=y\quad(y\in D(H)).
\]
For each output, its graph value is $x\mp iR R_\pm(R)x$. Finally require a joint PR approximation of each finite rational vector by graph members to any prescribed Hilbert error. Applying this to a finite sample supplies the same service for all names. These are first-order contracts on fixed presentations as in Section~\ref{sec:arithmetic}. One may equivalently describe a domain member by a range witness for $R_+(1)$. The equations imply $\|R_\pm(R)\|\le1/R$, by expanding $\|(H\pm iR)y\|^2=\|Hy\|^2+R^2\|y\|^2$. Symmetry gives $R_+^*=R_-$ by testing on their ranges. The graph is closed since $(x_n,Hx_n)\to(x,y)$ implies $x=R_+(1)(y+ix)$, which gives the requisite witness. Each use of convergence here means convergence with given PR moduli.

For clarity this is self-adjointness in the graph sense. If $w,z$ satisfy $\langle Hx,w\rangle=\langle x,z\rangle$ on the domain, then $\langle (H+i)x,w\rangle=\langle x,z-iw\rangle$. Substituting $x=R_+(1)v$ gives $w=R_-(1)(z-iw)$, hence $Hw=z$. Symmetry supplies the reverse inclusion. No Riesz representation principle producing an arbitrary functional's vector is required.

For a bounded self-adjoint $C$, define $e^{-itC}$ by its power series. With $|t|\|C\|\le M$ an integer, the norm tail after degree $n\ge 2M$ is at most $2^{1+4M-n}$: beyond $2M$ the ratio is at most $1/2$, and the starting term is at most $e^M\le 4^M$. For this scalar bound, $\sum 1/j!\le3$ by a geometric tail, and the finite multinomial expansion dominates any finite partial sum at integer $M$ by the $M$th power of that series. For $M=0$ all positive-degree terms vanish. Multiplication of absolutely convergent series is justified by taking rectangular truncations with the displayed tails; the binomial formula identifies the product at $t,s$ with the value at $t+s$. Taking adjoints gives unitarity. The same calculation applies to commuting bounded generators. Taylor expansion at zero and the group law give the derivative $-iCe^{-itC}$. Two useful inequalities are
\begin{equation}
 \|(e^{-itC}-I)v\|\le |t|\|Cv\|,\qquad
 \|(e^{-itC}-I+itC)v\|\le t^2\|C^2v\|/2.\label{eq:tay}
\end{equation}
One proof integrates the derivative, then integrates again; the derivatives are uniformly continuous on a compact interval with a PR modulus from the series. For an explicit justification of integration, approximate a Lipschitz vector function on an interval of length $T$ by a left sum of mesh $h$; comparison on a common refinement costs at most $T\,\mathrm{Lip}\, h$. A uniform modulus gives $T\omega(h)$ in the same way. For the fundamental theorem, telescope the local differences under a uniform derivative remainder $\|f(s+h)-f(s)-hf'(s)\|\le |h|\omega(|h|)$; its accumulated norm is at most $T\omega(h)$. Polynomial Taylor remainders on the group have such a uniform bound from unitarity and the series. Integrating the norm bounds $\|Cv\|$ and $s\|C^2v\|$ gives \eqref{eq:tay}. This argument also supplies the scalar exponential and its Lipschitz estimates used below.

\begin{theorem}[Resolvent data give unitary dynamics]\label{thm:stone-forward}
From the graph data above one constructs a strongly continuous unitary group $U(t)$ with joint PR evaluation on the presented inputs. On bounded time intervals each orbit has a PR modulus. Its strong generator has graph exactly $-iH$, with derivative witnesses on graph members. Conversely a named vector with a certified strong derivative is recovered as a graph member.
\end{theorem}
\begin{proof}
All resolvents commute. Indeed if $R_a=(H+a)^{-1}$ and $a,b$ are permitted imaginary shifts, then evaluating on $R_bx$ gives $R_a-R_b=(b-a)R_aR_b$; exchanging $a,b$ gives the second ordering, with identical shifts trivial. They commute with $H$ on its domain since $HR_a x=x-aR_a x$. Define the bounded self-adjoint operators
\[
 D_R=\frac R{2i}(R_-(R)-R_+(R)),\qquad H_R=\frac{R^2}{2}(R_+(R)+R_-(R)).
\]
The difference identity gives $D_R=R^2 R_+(R)R_-(R)$, whose range lies in $D(H^2)$: if $x$ is in the domain then $HR_a x=R_a Hx$ lies in the domain as well. Consequently $HD_R=H_R$, $H^2D_R=R^2(I-D_R)$, by applying the resolvent equations. We have $\|D_R\|\le1$, $\|H_R\|\le R$. On the domain,
\begin{equation}
 \|(I-D_R)y\|\le\|Hy\|/R,\qquad
 \|(H-H_R)v\|\le\|H^2v\|/R\quad(v\in D(H^2)).\label{eq:yosi}
\end{equation}
Indeed $I-D_R=R^{-2}H_R H$ there, by the preceding identity and commutation. For arbitrary $x$, approximate by a graph member $y$ at error $\epsilon$, obtaining $\|(I-D_R)x\|\le 2\epsilon+\|Hy\|/R$. A bound for $Hy$ is part of its name. This proves $D_Rx\to x$ strongly with an explicit schedule. Thus $D_Sx$ gives a $D(H^2)$ approximant to $x$, with $\|H^2D_Sx\|\le2S^2\|x\|$. For a graph input it also converges in graph norm since $HD_S x=D_S Hx$.

Put $U_R(t)=e^{-itH_R}$. The generators commute, so \eqref{eq:tay} applied to $H_R-H_S$ gives
\[
 \|(U_R(t)-U_S(t))v\|\le |t|(R^{-1}+S^{-1})\|H^2v\|\qquad(v\in D(H^2)).
\]
Given $x$, $|t|\le T$, and $\epsilon=2^{-k-5}$, obtain $v=D_bx$ within $\epsilon$ by the preceding rule. Let $K\ge 2b^2\|x\|$ be an integer. For $R,S\ge1+\lceil 2TK/\epsilon\rceil$, the full-vector difference is at most $3\epsilon$. This constructs $U(t)x$ from a finite polynomial computation and a diagonal precision request. The group law survives since all approximate maps are isometries: to pass two factors to the limit compare first on $x$, then on the fixed named vector $U(s)x$. The inverse is $U(-t)$ and the adjoint is its inverse. For strong continuity, take a graph approximant $y$ to $x$ within $\epsilon$ and use \eqref{eq:tay} and $H_R y=D_R Hy$ to get $\|(U(h)-I)x\|\le2\epsilon+|h|\|Hy\|$. The group law makes this modulus uniform in the base time.

All $U(t)$ commute with $R_+(1)$ by their construction; hence they preserve the graph and commute with $H$ there. On $D(H^2)$, \eqref{eq:tay} and $H_R^2v=D_R^2H^2v$ give
\[
 \|U(h)v-v+ihHv\|\le h^2\|H^2v\|/2 .
\]
For general $x\in D(H)$, approximate $x$ in graph norm by $v=D_bx$ with $\|Hx-Hv\|\le\epsilon$. The first bound in \eqref{eq:tay} passes to the limit on graph inputs. The difference quotient error for $x$ is therefore at most $2\epsilon+|h|\|H^2v\|/2$ for $h\ne0$. This is a PR derivative modulus. Conversely, suppose the difference quotient on $x$ tends in norm to a named $z$ with a given modulus. Applying $R_+(1)$ to it and using the already proved derivative on $R_+(1)x$ gives $R_+(1)z=-iH R_+(1)x$. Rearrangement gives $x=R_+(1)(iz+ix)$, so $Hx=iz$. All quotients suffice at rational $h\ne0$; the same estimates extend to named $h$ with an apartness witness.
\end{proof}

\begin{theorem}[A named unitary group supplies resolvents]\label{thm:stone-back}
Suppose that a unitary group is presented by jointly PR approximations on rational times and finite rational inputs, with the group laws and a single joint PR strong continuity modulus $m(q,k)$ at zero for the finite input $q$ and precision $k$. It extends to named vectors and times. Its certified strong generator determines self-adjoint graph data as above. The two constructions return the original group and graph extensionally.
\end{theorem}
\begin{proof}
Approximate an arbitrary vector $x$ by a finite sample $q$ within $\epsilon$. For $h$ within the supplied neighborhood for $q$, the displacement of $x$ is at most $2\epsilon+\|(U(h)-I)q\|$. Unitarity transports this estimate to every time. Approximating the time by a rational now extends the orbit. Initially use a witnessed relation: for each supplied PR modulus define $Hx=iz$ when $z$ is a named norm derivative at zero with a PR modulus; the first-order law tests rational nonzero increments. The difference quotient is linear and the derivative is unique by removing arbitrary slack. Differentiating $\langle U(h)x,U(h)y\rangle=\langle x,y\rangle$ on derivative inputs gives symmetry, since the first factor is conjugate-linear. Expansion at finite $h$ bounds the product remainder by the two supplied remainder bounds and a term proportional to $|h|$.

For $R\ge1$ set
\[
 R_-(R)x=i\int_0^\infty e^{-Rs}U(s)x\,ds,\qquad
 R_+(R)x=-i\int_0^\infty e^{-Rs}U(-s)x\,ds .
\]
For $\|x\|\le M$ the tail beyond integer $b$ is at most $Me^{-Rb}/R\le M2^{-b}$ (use $e\ge2$). On $[0,b]$ the integrand has modulus from that of the orbit plus $RM|s-s'|$ since the scalar damping has derivative of magnitude at most $R$. The elementary Riemann-sum rule above gives joint PR samplers. If $r$ denotes the first integral without $i$, changing variables gives, for signed rational $h$,
\[
 U(h)r=e^{Rh}\left(r-\int_0^h e^{-Rs}U(s)x\,ds\right).
\]
Translations of Riemann grids prove the formula first on finite intervals, then the exponential tail proves it at infinity. On $|h|\le1/R$, the average integrand over the oriented short interval differs from $x$ by at most $3\sup_{|s|\le|h|}\|(U(s)-I)x\|+3RM|h|$. Also $|(e^{Rh}-1)/h-R|\le2R^2|h|$ by Taylor's formula. Thus $r$ has derivative $Rr-x$ with a PR modulus, and $(H-iR)i r=x$. Apply the same argument to $U(-t)$ for the other sign. Symmetry gives the displayed norm identity at the beginning of this section, proving injectivity and the left inverse on the domain. Adjointness follows as before. The resulting graph has the single test $R_+(1)(y+ix)=x$. The integral calculation supplies a derivative modulus for every pair passing this test, and the inverse identity proves that every witnessed derivative pair passes it. This presentation removes the apparent quantifier over moduli.

To approximate $x$ in the graph domain, use $Rr$, since $\int_0^\infty Re^{-Rs}\,ds=1$. Choose rational $a>0$ so the orbit displacement on $[0,a]$ is at most $\epsilon$. Splitting the integral there bounds $\|Rr-x\|$ by $\epsilon+2M e^{-Ra}\le\epsilon+2M/(1+Ra)$. Take $R\ge1+\lceil2M/(a\epsilon)\rceil$. This supplies the required graph approximant with its derivative data.

Finally the resolvents integrate the group itself, so commute with that group. The reconstructed group $W(t)$ therefore also commutes with $U(s)$. On a graph vector $x$ both have derivative $-iHx$ at zero. For a fixed rational $t$ the commuting-power telescope bounds $\|(U(t)-W(t))x\|$ by $n\|(U(t/n)-W(t/n))x\|$. For any target error, choose $n$ by the two derivative moduli, making the latter product arbitrarily small. Density extends the identity to all vectors, then continuity extends it to every named time. The converse graph identification is already part of Theorem~\ref{thm:stone-forward}.
\end{proof}

\subsection{Driven systems and the two pictures of dynamics}
For a time-dependent bounded self-adjoint $K(t)$ on a compact time window, assume a joint PR column sampler, an integer bound $B$ and a PR \emph{operator-norm} continuity modulus. Define a propagator by freezing $K$ at the left endpoint of each interval and multiplying the unitary exponentials, with later factors on the left. This prescription describes a spin in a changing magnetic field, including noncommuting field directions.

For bounded self-adjoint $A,B'$ the path $e^{-i(h-u)A}e^{-iuB'}$ has derivative $ie^{-i(h-u)A}(A-B')e^{-iuB'}$. Integration as in \eqref{eq:tay} gives $\|e^{-ihA}-e^{-ihB'}\|\le |h|\|A-B'\|$, without a commutation assumption. Let the total time length be at most $T$. Refining any partition of mesh $d$ changes the product by at most $T\omega(d)$, where $\omega(d)$ bounds the variation in $K$ over any time span $d$. Indeed split each frozen exponential exactly on the refined partition, then telescope the local changes; the other factors all have norm one. Two partitions have a common refinement, so the products converge in the sense of a uniform operator-error bound. For a tolerance $\delta_k$, choose a dyadic mesh on which $\omega(d)\le\delta_{k+3}/(1+T)$. Variable products of exponentials on finite-vector samplers use the accumulation rule of Section~\ref{sec:arithmetic}. Reversing each product gives the limiting inverse and adjoint. Concatenation, via partitions containing the intermediate time, proves $U(t,r)U(r,s)=U(t,s)$; the time endpoints extend to names since the displacement over an interval of length $h$ is at most $B|h|$ on unit vectors.

On a short interval starting at $t$, the freezing estimate and \eqref{eq:tay} bound the error relative to $I-ihK(t)$ by $|h|\omega(|h|)+h^2B^2/2$, for either sign (for a backwards interval use the reversed product). Thus $\partial_t U(t,s)x=-iK(t)U(t,s)x$ with a computable modulus. In finite dimension set $\rho(t)=U(t,s)\rho(s)U(t,s)^*$ and $A_H(t)=U(t,s)^*A(t)U(t,s)$. The finite trace identity gives $\operatorname{tr}(\rho(t)A(t))=\operatorname{tr}(\rho(s)A_H(t))$. The product difference formula and the derivative estimates yield
\[
 i\,\dot\rho=[K,\rho],\qquad
 \dot A_H=U^*\bigl(\dot A+i[K,A]\bigr)U .
\]
For a time-dependent observable assume a PR derivative modulus; for a constant $A$ commuting with $K(t)$ its expectation is conserved. For bounded tests in a purified state the identical calculation pairs vectors instead of tracing matrices. If energy is unbounded the product rule on a vector requires the relevant graph inputs; Section~\ref{sec:dynamics} provides a common invariant domain for the oscillator identities.

\section{Energy distributions for general presented Hamiltonians}\label{sec:unbounded-calculus}
The resolvent data of Section~\ref{sec:stone} also suffice to model a detector whose bounded continuous response to energy is known only on compact intervals. This matters when the energy distribution has no global support bound and no finite second moment in the prepared state.

\begin{lemma}[Bounded transform with a dense range witness]
Put $D=D_1$, $C=H_1=HD$ and $S=\sqrt D$. There is a presented self-adjoint contraction $F$ with $FS=SF=C$ and $I-F^2=D$. For every named $x$ one can approximate $x$ to any requested precision by $Dy$ with an explicit name and bound for $y$.
\end{lemma}
\begin{proof}
$D=R_+(1)R_-(1)$ is positive and $D\le I$. The resolvent identities give $C^2=D(I-D)$, first on the range construction, or directly by multiplying $H^2D=I-D$ by $D$. Thus $\|Cz\|^2\le\|Sz\|^2$ for every $z$. Given an approximant $v=D_bx$ from Theorem~\ref{thm:stone-forward}, set $y=v+H^2v$. Then $Dy=v$, by the inverse equations for $H+i$ and $H-i$ on $D(H^2)$, and $\|y\|\le(1+2b^2)\|x\|$. This supplies the claimed witness. Since $Dy=S(Sy)$ it also supplies dense range approximants for $S$.

On that range define $F(Sz)=Cz$. The estimate just proved makes this assignment well-defined and contractive. At precision $k$, choose $v=Dy$ within $\delta_{k+3}$ of the input, and compute $C S y$ within $\delta_{k+3}$; comparing approximants in the range by contraction gives a regular output. This bounded extension is self-adjoint because $C$ and $S$ commute and $\langle Cz,Sw\rangle=\langle Sz,Cw\rangle$. Taking adjoints of $FS=C$ gives $SF=C$. Thus $(F^2-I+D)D=0$ by $C^2=D(I-D)$, and the range approximation just proved yields $F^2=I-D$ on the entire space. Classically $F$ represents $H/\sqrt{1+H^2}$. The algorithm never applies an inverse square root at zero.
\end{proof}

\begin{theorem}[Bounded continuous energy responses]\label{thm:unbounded-calculus}
Let $f:\mathbb R\to\mathbb C$ be presented by a global integer bound $M\ge1$ and joint PR sampling and continuity moduli on rational compact intervals. For the resolvent data there is an extensional positive, unital, complex-linear, multiplicative, adjoint-preserving assignment $f\mapsto f(H)$ with $\|f(H)\|\le M$. It gives the supplied resolvents at $\pm i$ and the unitary evolution for the phase response $f(\lambda)=e^{-it\lambda}$. A normalized state's expectations have effectively tight energy tails in the continuous-detector sense below.
\end{theorem}
\begin{proof}
For $-1\le s\le1$, put $d(s)=1-s^2$. For dyadic $0<a\le1/2$ define
\[
 h_a(s)=\max(0,\min(1,(d(s)-a)/a)),\qquad
 g_a(f)(s)=h_a(s) f\!\left(\frac{s}{\sqrt{\max(a,d(s))}}\right).
\]
These are presented continuous functions on a fixed compact interval. For instance the argument of $f$ has range in $[-1/a,1/a]$ and Lipschitz bound $1/\sqrt a+1/a^{3/2}$, by the difference of inverse square roots above $a$ and the bound $|d(s)-d(s')|\le2|s-s'|$. The sampler on rationals extends to such named arguments by the given local modulus. Section~\ref{sec:spectrum} evaluates $g_a(f)(F)$ by rescaling the interval.
For $0<b\le a$ both cutoffs agree with the uncropped value where $d(s)\ge2a$, and everywhere $|g_a(f)-g_b(f)|\le2M$. It follows from the polynomial/continuous norm inequality that
\[
 \|(g_a(f)(F)-g_b(f)(F))D\|\le4Ma.
\]
For an arbitrary $x$ choose $Dy$ within $\epsilon/(4M)$ by the lemma and then $a\le \epsilon/(8M(1+\|y\|_{\rm bound}))$. The two actions on $x$ differ by at most $\epsilon$. Their strong limit defines $f(H)$ with the claimed bound. This proof supplies a joint PR algorithm even when the cutoffs fail to converge in operator norm. Conjugating the scalar function gives an adjoint on the dense tests. Linearity is immediate. For two functions $f_1,f_2$, the discrepancy of $g_a(f_1)g_a(f_2)$ from $g_a(f_1 f_2)$ vanishes where $d\ge2a$ and is bounded by $2M_1M_2$. The same dense-range estimate sends the product discrepancy strongly to zero, and uniformly bounded strong limits pass through products by insertion of a fixed vector. Applying this argument to $h_a-1$ gives unitality. Nonnegative $f$ gives positive approximants and positive limit.

For $f(\lambda)=(\lambda\pm i)^{-1}$, the uncropped function extends continuously to $s=\pm1$ by $r_\pm(s)=s\sqrt{d(s)}\mp i d(s)$. The difference from the cutoffs vanishes where $d\ge2a$ and has a fixed bound, hence $f(H)=FS\mp iD=C\mp iD=R_\pm(1)$. The same argument identifies any compactified function with a continuous extension in $s$. For the phase response write $V(t)=f_t(H)$. Multiplicativity and conjugation make $V$ a unitary group. For $x=Dy$ the scalar Taylor inequality on the range $d>0$ gives
\[
 \|V(t)Dy-Dy+itCy\|\le \tfrac12 t^2\|y\|.
\]
Indeed multiply the scalar remainder by $d$: its quadratic bound is $t^2s^2/2$. The functions $d(s)e^{-its/\sqrt d}$ and $s\sqrt d$ extend continuously by zero at the endpoints, with effective moduli by cutting off where $d\le2a$ (respective errors at most $2a$ and $\sqrt{2a}$). In particular $d\,g_a(f_t)$ approximates the first extended function uniformly within $4a$; thus its compact calculus is exactly $V(t)D$. The scalar Taylor bound now proves the estimate including the endpoints. The group $U$ of Theorem~\ref{thm:stone-forward} commutes with $D,C,S$ and with $F$, the last by $FS=C$ and the effective dense range of $S$. It therefore commutes with $V$. Their individual quadratic remainder estimates on $Dy$ give $\|(V(h)-U(h))Dy\|\le h^2(\|y\|+\|(I-D)y\|)/2\le h^2\|y\|$. By commuting and telescoping $n$ factors with $h=t/n$, the difference at $t$ is at most $t^2\|y\|/n$. Density extends equality to all $x$.

For a unit vector $x$, compute $Dy$ close to $x$. The effect $I-h_a(F)$ has expectation at most $\|x-Dy\|+2a\|y\|$, since $0\le1-h_a\le1$ and $(1-h_a)d\le2a$. Under the energy parametrization $d=1/(1+\lambda^2)$, the detector $h_a$ vanishes for $|\lambda|\ge1/a$ and equals one when $d\ge2a$. Thus one can place all but any prescribed weight in a bounded energy region measured with a continuous cutoff. For a purified mixed state apply the same proof on the system tensored with the identity on the ancilla. The range of $D\otimes I$ has joint PR dense approximants by resolving the finitely many slices of a finite array and assigning the requested error divided by one plus the number of slices. This establishes tightness without an energy-expectation input. Ramps and boundary-mass witnesses give individual threshold probabilities as in Section~\ref{sec:spectrum}.
\end{proof}

\section{Hamiltonians with quantitative unitary dynamics}
\label{sec:dynamics}

We use units in which $\hbar=1$. An unbounded observable is specified by its
\emph{graph}: a state vector $x$ in its domain comes with a second Hilbert
name $y$ and witnesses of the equations that say $y=Hx$. The name of $y$
supplies information that can be stronger than a bound on $\|y\|$.
For example, differentiating a trajectory in Hilbert norm requires the
tail of its derivative to be small at a known cutoff.

A strong derivative certificate for a curve $X$ at $t$ consists of a
Hilbert name $z$ and a PR function $d(k)$ such that
\[
 \|X(t+h)-X(t)-hz\|\le 2^{-k}|h|
 \quad\text{whenever } |h|\le 2^{-d(k)} .
\]
One may test rational increments and extend to named increments by
the continuity estimates, using rational approximants inside the closed
interval. To compare two proposed derivatives, substitute the positive
rational $h=2^{-\max(d_1(k),d_2(k))}$; rational cancellation gives
$\|z_1-z_2\|\le2^{1-k}$ for arbitrary $k$.

\subsection{Energy coordinates}

Let $e_j$ be the canonical basis of $\ell^2(\mathbb N)$. Supply real
energy names $\lambda_j$ uniformly in $j$, with a PR upper bound
$\Lambda(N)\ge 1+\max_{j<N}|\lambda_j|$; an integer bound of this kind
can always be computed from the stage-zero approximations. Define
\[
 (x,y)\in D \ \Longleftrightarrow\ (\forall j)\ y_j=\lambda_jx_j,
 \qquad (U(t)x)_j=e^{-it\lambda_j}x_j .
\]
Here and below a universal statement about the coordinates abbreviates
tests at arbitrary finite index and precision. Write $P_N$ for the first
$N$ coordinate projection and $Q_N=I-P_N$. A finite approximation $v$
to $x$ within $\epsilon$, supported below $N$, proves
$\|Q_Nx\|\le\epsilon$ by contraction.

\begin{theorem}[Diagonal graph and generator]\label{thm:diag-energy}
$D$ is a densely defined, closed, self-adjoint graph. The maps $U(t)$
form a strongly continuous unitary group whose strong derivative at zero
has graph $-iD$. For every graph pair $(x,y)$, $U(t)x$ has derivative
$-iU(t)y$ at $t$. All constructions have PR realizers.
\end{theorem}
\begin{proof}
Uniqueness and closedness follow by applying a coordinate functional to
the approximants: equal coordinates give equal vectors by truncating both
tails. Every basis vector is in the domain, and the graph truncations
$(P_Nx,P_Ny)$ converge in both components with an explicit joint cutoff.
The adjoint graph consists of pairs $(x,z)$ with
$\langle Dv,x\rangle=\langle v,z\rangle$ on graph inputs $v$
(inner products are linear in the second slot).
Testing $v=e_j$ forces $z_j=\lambda_jx_j$.
Conversely the equalities on finite sums pass to graph inputs by the joint
cutoff and Cauchy--Schwarz. The resolvents $(D\pm i)^{-1}$ are diagonal
multipliers by $(\lambda_j\pm i)^{-1}$. Both these coefficients and
$\lambda_j/(\lambda_j\pm i)$ have modulus at most one. This supplies
names for the solution and its graph partner for every right-hand side.

For computation of a phase or a finite matrix exponential, the elementary
bound
\begin{equation}\label{eq:exp-pr}
 \sum_{r=m}^\infty {a^r\over r!}
 \le 2^{-m}\sum_{r=0}^\infty{(2a)^r\over r!}
 \le 2^{4L-m}\qquad(0\le a\le L,\ L\in\mathbb N)
\end{equation}
follows by multiplying the series with positive coefficients and using
$\sum 1/r!\le3$, hence $\exp(2L)\le3^{2L}$. These series identities can
equally be verified by the binomial convolution of finite coefficients
and the geometric tail once successive terms have ratio at most $1/2$.
Absolute product and adjoint identities give
$e^{is}e^{it}=e^{i(s+t)}$ and $|e^{it}|=1$ for real $s,t$.
Integrating the convergent series on a bounded interval gives
\[
 |e^{is}-1|\le |s|,\qquad |e^{is}-1-is|\le s^2/2 .
\]
For error $\epsilon=2^{-k}$, take a finite $v$ within $\epsilon/4$ of
$x$. On its support use uniform phase errors at most
$\epsilon/(4(1+\|v\|))$ and form a finite output to accuracy
$\epsilon/4$. The parameter $L$ in \eqref{eq:exp-pr} is bounded from
$|t|$ and $\Lambda(N)$. For
$|t|\le T$ and $|\lambda|\le M$, rational approximations to $t,\lambda$
within $\rho\le1$ give a product error at most $(T+M+1)\rho$;
the unit-phase map is $1$-Lipschitz. All norms appearing in precision choices may be
rounded upwards using rational upper estimates. Coordinate multiplication
proves the group and adjoint laws. On a head and tail one obtains
\[
 \|(U(h)-I)x\|\le |h|\Lambda(N)\|x\|+2\|Q_Nx\|.
\]
This is a strong continuity modulus at every $t$. For a graph pair the
first-order remainder satisfies
\[
 \|U(h)x-x+ihy\|
 \le |h|\bigl(2\|Q_Ny\|+|h|\Lambda(N)^2\|x\|/2\bigr).
\]
Indeed, on each tail coordinate the first-order difference is bounded
by $2|h||y_j|$, and on the head the second scalar bound applies.
Choose $N$ from the name of $y$ and then a sufficiently small dyadic
bound on $|h|$. Conversely, a named derivative $z$ implies, by the same
scalar expansion at each fixed $j$, that $z_j=-i\lambda_jx_j$.
Thus $(x,iz)$ is a graph pair. Translating time gives the remaining
claim.
\end{proof}

\subsection{A number-basis oscillator}

On finite sequences put $ae_0=0$,
$ae_n=\sqrt n\,e_{n-1}$ for $n>0$, and
$a^\dagger e_n=\sqrt{n+1}\,e_{n+1}$. The square roots are obtained
by bounded bisection. Set
\[
 Q=(a+a^\dagger)/\sqrt2,\qquad
 P=(a-a^\dagger)/(i\sqrt2),\qquad
 H=(P^2+Q^2)/2 .
\]
The core of finite sequences is invariant under every displayed
operator. Inner products of basis vectors give the stated adjunction.
Acting on $e_n$ gives $aa^\dagger e_n=(n+1)e_n$ and
$a^\dagger a e_n=ne_n$, including $n=0$. Expanding products therefore
proves on the core
\[
 [Q,P]=iI,\qquad H=a^\dagger a+\tfrac12 I,\qquad
 U(t)e_n=e^{-it(n+1/2)}e_n .
\]
The graph closure of this core Hamiltonian is exactly the graph of
Theorem~\ref{thm:diag-energy} with $\lambda_n=n+1/2$, by joint
truncation. Energy measurement in the number basis gives outcome $n+1/2$
with probability $|x_n|^2$; omitted mass equals $\|Q_Nx\|^2$.
On the core $U(t)^*aU(t)=e^{-it}a$ and
$U(t)^*a^\dagger U(t)=e^{it}a^\dagger$. Recombining,
$Q(t)=Q\cos t+P\sin t$ and $P(t)=P\cos t-Q\sin t$.
For unit core input $x$, take the centered vectors
$u=(Q-\langle x,Qx\rangle)x$, $v=(P-\langle x,Px\rangle)x$.
The inner-product difference
$\langle u,v\rangle-\langle v,u\rangle=i$ follows from the commutator.
Cauchy--Schwarz gives $\Delta Q\,\Delta P=\|u\|\|v\|\ge1/2$.
At $x=e_0$ the centered vectors are $e_1/\sqrt2$ and $ie_1/\sqrt2$.
Hence equality holds. The commutator identity has been asserted on
a specified common domain throughout.

\subsection{Perturbing an unbounded Hamiltonian}

The next construction gives a state-dependent spatial cutoff for a
Hamiltonian that need not be diagonal in the given basis. It can be used
when the resolvent of the interacting operator has not been computed.

\begin{theorem}[Finite-band perturbation]\label{thm:band-dynamics}
Let $V$ have uniformly named entries, $V_{jk}=\overline{V_{kj}}$,
$V_{jk}=0$ for $|j-k|>w$, and a supplied integer bound $B\ge\|V\|$
on finite vectors. On the graph domain of $D$, set $H=D+V$.
This graph is self-adjoint and generates a PR-presented unitary group
$W(t)$. If $v$ is supported below $J$, $|t|\le T$ for an integer $T$,
and $K\ge J+mw$ with $m\ge1$, then
\begin{equation}\label{eq:galerkin}
 \|W(t)v-e^{-itP_KHP_K}v\|
 \le 2\|v\|\,(BT)^m/m!\le2\|v\|\,2^{4BT-m}.
\end{equation}
\end{theorem}
\begin{proof}
A finite input produces finite output under $V$; its norm bound extends
$V$ and its adjoint to all names. The adjoint relation for $H$ is the
one for $D$ after moving the bounded $V$ to the other side. Closedness
follows by subtracting $Vx$ from the second graph component. Joint
truncations for $D$ approximate $H$ in graph norm as well.

Write $D_K=P_KDP_K$, $V_K=P_KVP_K$, $H_K=D_K+V_K$.
The finite matrix exponential, extended by the identity on the
complement, is evaluated by \eqref{eq:exp-pr} using
$L\ge T(\Lambda(K)+B)$. Conjugate transpose and the group law follow by
absolutely convergent binomial convolution for commuting powers.

Here is a finite-product proof of the cutoff estimate. For $n\ge1$
put $h=t/n$, $A_K=e^{-ihD_K}$, $C_K=-ihA_KV_K$, and $F_K=A_K+C_K$.
For any finite self-adjoint $L$,
$\|e^{-ihL}-I\|\le |h|\|L\|$ and
$\|e^{-ihL}-I+ihL\|\le h^2\|L\|^2/2$ by integrating its unitary
series once and twice. Expansion to first order consequently gives
\[
 \|e^{-ihH_K}-F_K\|\le h^2(\Lambda(K)+B)^2,\qquad
 \|F_K\|^2\le1+h^2 B^2 .
\]
For the second bound expand $\|v-ihV_Kv\|^2$; the mixed term vanishes.
Telescoping $n$ factors yields
\[
 \|e^{-itH_K}-F_K^n\|
 \le T^2(\Lambda(K)+B)^2 e^{T^2 B^2/(2n)}/n.
\]
Every limit in $n$ at a \emph{fixed} cutoff therefore has a PR rate.

Expand $F_K^n$ from the left, choosing $A_K$ or $C_K$ at each
factor, stopping a branch immediately on choosing its $m$th $C_K$.
The complete branches with $r<m$ kicks apply $V_K$ at most $r$ times
to $v$. Between kicks there are diagonal phases. These branches agree
for all $K\ge J+mw$, since each kick increases the support by at
most $w$. Each stopped branch has $m$ kicks, diagonal phases, and a
rightmost block of unexpanded $F_K$'s. Its norm on $v$ is at most
$(|h|B)^m e^{T^2B^2/(2n)}\|v\|$. There are $\binom nm$ such branches
(with zero branches for $m>n$). Since
$\binom nm/n^m\le1/m!$, subtracting the two cutoff expressions
costs at most $2\|v\|(BT)^m e^{T^2B^2/(2n)}/m!$.
For any \emph{two} specified cutoffs send $n$ to infinity using the preceding
telescoping estimate with the larger of the two energy bounds. This proves the
pairwise version of \eqref{eq:galerkin}. The expansion is a finite
subset enumeration with at most $2^n$ full words. One may replace
the finite matrices by their rational approximants throughout, budgeting
a local error $\eta$ per factor by $n\eta(M+1)^n$ for input norm one
and factor bounds $M$ when $\eta\le1$. These estimates give the
finite-precision induction directly.

Given arbitrary $x$ and error $\epsilon$, approximate by finite $v$
within $\epsilon/8$, take
$m\ge k+4BT+4+C$ with an integer $C\ge\|v\|$ for $\epsilon=2^{-k}$, and evaluate
the chosen matrix exponential at $v$ within $\epsilon/8$.
The Cauchy estimate constructs the strong limit, uniformly in bounded
time, and proves \eqref{eq:galerkin}. It preserves norms. Products pass
to the strong limit: insert the fixed vector $W(s)x$ between the
two approximants and use their contraction bounds. The finite
composition law gives $W(t)W(s)=W(t+s)$, and $W(-t)$ is its inverse.
Strong continuity follows from approximation by one finite evolution
uniformly on a time interval.

For a finite $v$, all sufficiently large cutoffs agree on $H v$ and
$H^2v$. Twice integrating the finite unitary evolution gives
\[
 \|(W(h)-I)v\|\le |h|\|Hv\|,\qquad
 \|W(h)v-v+ihHv\|\le |h|^2\|H^2v\|/2.
\]
The first inequality extends to every graph input by joint truncation.
For such an $x$ the derivative remainder is bounded by $|h|$ times $2\|H(x-P_Nx)\|+|h|\|H^2P_Nx\|/2$. Its first term is
bounded by $2(\|Q_NDx\|+B\|Q_Nx\|)$.
Conversely, if the trajectory at $x$ has derivative $z$, differentiate
$\langle e_j,W(h)x\rangle=\langle W(-h)e_j,x\rangle$ at zero.
This gives $(iz)_j=\lambda_j x_j+(Vx)_j$. Hence $iz-Vx$ is the
required name of $Dx$, identifying the entire generator graph.
The group law now gives the differential equation on this domain
at every time, including invariance of the domain.
\end{proof}

For instance $V=g(S+S^*)$ for the unilateral shift $Se_j=e_{j+1}$
and real coupling $g$ has $w=1$, $B\ge2|g|$. Taking
$\lambda_j=(j+1/2)^2$ gives an unbounded anharmonic level sequence
with adjacent-level mixing. The cutoff estimate is independent of the
magnitude of the largest retained energy; that magnitude affects the
cost of the finite exponential.

\section{A particle on the line}
\label{sec:particle}

Continuous spectrum can be treated using the same arithmetic of finite
approximants. Take Gaussian-rational step functions of compact support on
the real line; a basic interval is $[u,v)$ with rational endpoints.
Refining to a common finite partition makes
\[
 \langle f,g\rangle=\sum_I |I|\,\overline{f_I}g_I
\]
a rational complex computation. The Hilbert space used here, denoted
$L^2(\mathbb R)$, is the regular metric completion of these step functions.
Values of a general vector at a point have no role in the definition.
One can identify this presentation with the coordinate Hilbert space
explicitly. Replacing endpoints by dyadics within $\delta$ costs at most
$\sqrt{2\delta}\sum_I|f_I|$ by the triangle inequality. On each integer
unit interval take its normalized indicator and the Haar differences
$|J|^{-1/2}(1_{J_{\rm left}}-1_{J_{\rm right}})$ for dyadic subintervals
$J$ of that unit interval. Orthogonality follows by cancellation on the
smaller interval. On each finite dyadic partition the indicator of either
half is a linear combination of the parent indicator and its difference.
Successive subdivision therefore transforms every compact dyadic step
into a finite Haar expansion preserving its squared norm, with finite
inverse. Integer intervals and their dyadic nodes admit a PR enumeration
with PR inverse (encode the signed integer and the depth/offset by pairing).
The finite transforms and density estimates extend both isometries to
the regular completions.

\subsection{Momentum, translations and the free Hamiltonian}

Let $R$ be a positive integer and write $E_R$ for multiplication by
$1_{[-R,R)}$. It is an orthogonal projection of norm at most one.
Every named vector has a PR cutoff for $\|(I-E_R)x\|$ by taking the
radius of a sufficiently accurate step approximation.
For a real multiplier $b(p)$ with supplied PR bounds $M_R$ and
continuity moduli on $[-R,R]$, bounded multiplication by $E_Rb$
is computed on steps by subdivision. Local oscillation $\eta$ costs
at most $\eta\|f\|$ on a step input $f$; the bounded extension
costs at most $M_R$ times its input error. Define the graph
\[
 (x,y)\in H_b \quad\Longleftrightarrow\quad
 E_Ry=(E_Rb)x\ \text{ for every positive integer }R.
\]
Suppose also that $b$ has uniformly PR point approximants on each such
interval. This graph is closed by boundedness of the local multipliers.
Compact steps lie in the domain. The graph pairs
$(E_Rx,E_Ry)$ approximate any graph pair, with both tails supplied
by the names. To identify the adjoint, test its identity against
arbitrary compact steps supported on $[-R,R)$. The bounded real
multiplier there is self-adjoint by the finite integral and passage to
the limit. Density of those steps gives exactly $E_Ry=(E_Rb)x$.
Conversely the identity on bounded support extends by truncating both
graph components. Thus $H_b$ is self-adjoint.

Multiplication by $e^{-itb(p)}$ yields a strongly continuous unitary
group. On a fixed compact interval the modulus of this multiplier follows
from $|e^{is}-e^{it}|\le |s-t|$; the input support controls the tail.
The two head--tail estimates of Theorem~\ref{thm:diag-energy} hold
with $P_N$ replaced by $E_R$, $\Lambda$ by $M_R$, and $Q_N$ by $I-E_R$.
Conversely, projecting a putative derivative onto $[-R,R)$ identifies
it with $-i(E_Rb)x$ using the bounded multiplier Taylor estimate.
This proves the strong Schr\"odinger equation with precisely the graph
domain displayed above.

For the momentum observable choose $b(p)=p$, and for the free
Hamiltonian of mass one choose $b(p)=p^2/2$. The latter group is
\[
  (U_0(t)f)(p)=e^{-itp^2/2}f(p).
\]
Momentum in a rational interval $I$ is measured by $f\mapsto1_I f$.
In a unit state the probability is $\|1_I f\|^2$. The energy interval
$[a,b]$, for rational $0\le a\le b$, is obtained by retaining the
two intervals with positive absolute momentum between $\sqrt{2a}$
and $\sqrt{2b}$. Endpoint inclusions represent the same vector.
These projections are effective: approximate both endpoints by rational
numbers within $\delta$, then use the step-height bound times the square
root of the total displaced length. For general input first approximate
by a step and use contraction. This supplies the sharp Born probabilities
even at an energy threshold.

\begin{proposition}[Continuous energy spectrum, with witnesses]
For each named $E\ge0$ and positive dyadic $\epsilon$ there is a
PR-prepared unit vector $v$ with $\|(H_{p^2/2}-E)v\|\le\epsilon$.
Every named solution of $H_{p^2/2}x=Ex$ equals zero, and
$H_{p^2/2}-E$ has dense proper range with explicit witnesses.
For $E\le-\eta$ with a supplied rational $\eta>0$ the bounded inverse
of $H_{p^2/2}-E$ exists and has bound $1/\eta$.
\end{proposition}
\begin{proof}
Put $s=\sqrt{2E}$ and take an integer $A\ge 2+s$.
Choose positive dyadic $\delta\le\min(1/4,\epsilon/(4A))$ and rational
$q$ with $|q-s|\le\delta$. Use the normalized indicator of
$[q-\delta,q+\delta)$. On its support $|p-s|\le2\delta$ and
$|p+s|\le2s+2\delta$, so its residual norm is at most
$2\delta(s+\delta)\le\epsilon$.

For a supposed eigenvector, approximate $x$ by a step $f$ within
$\gamma>0$ with height bound $C$. Use the union $J$ of the two intervals
of radius $2\delta$ about $q$ and $-q$, with $|q-s|\le\delta$.
Off $J$, both $|p-s|$ and $|p+s|$ are at least $\delta$.
On every compact truncation outside $J$, multiplication by
$|p^2/2-E|$ bounds the norm below by $\delta^2/2$, hence
the eigenvector equation forces that truncation to vanish. Passing
to the known spatial tail gives $x=1_Jx$. Its norm is bounded
by $\gamma+C\sqrt{8\delta}$. Taking
$\delta\le\gamma^2/(8(C+1)^2)$ proves $x=0$.

For density of the range, apply the same cut to a compact step $f$.
The vector $f1_{J^c}$ differs from $f$ by at most $C\sqrt{8\delta}$,
and dividing it by $p^2/2-E$ on its support gives a named graph
preimage. Indeed the reciprocal there is bounded by $2/\delta^2$;
on the finitely many retained compact intervals its Lipschitz bound
is at most $4D/\delta^4$ for support in $[-D,D]$, so it can be
sampled directly. Define its value as zero on the discarded intervals.
For properness let $g=1_{[s,s+1)}$, a named vector by endpoint
approximation. If $(H_{p^2/2}-E)x=g$, projection onto $[s,s+d)$
for rational $0<d\le1$ gives
$\sqrt d\le d(s+1)\|x\|$. This contradicts the choice of
$d$ with $\sqrt d(s+1)\|x\|\le1/2$, available from rational
upper bounds. Thus the asserted continuous-spectrum behavior has
finite-observation witnesses for density and for the obstruction to inversion.
For negative separated $E$, multiplication by $1/(p^2/2-E)$
has the stated bound and locally computable moduli; the graph partner
is the input plus $E$ times the output. The approximate eigenvectors
also rule out a supplied bounded inverse at any $E\ge0$.
\end{proof}

\subsection{Fourier transport with an explicit norm estimate}

Use $p$ for momentum and $x$ for position. For a finite momentum step on a disjoint partition
$f=\sum_j c_j1_{[a_j,b_j)}$ put
\[
 ({\cal F}f)(x)=(2\pi)^{-1/2}\int e^{ixp}f(p)\,dp,\quad
 S=\sum_j|c_j|,\quad A=\int |f|,\quad D_f=\int |p f(p)|\,dp .
\]
All displayed constants can be replaced by rational upper bounds.
The scalar constant $\pi$ is computable, for example from
$4\sum_{j=0}^{n-1}(-1)^j/(2j+1)$ with error at most $4/(2n+1)$.
The integral defining $\mathcal F f$ uses finite intervals. For $|x|\le R$
one may integrate the exponential Taylor polynomial term by term, with
\eqref{eq:exp-pr} and $L$ bounding $R\max_j(|a_j|,|b_j|)$.
The transform has height at most $A/\sqrt{2\pi}$ and Lipschitz
bound $D_f/\sqrt{2\pi}$. For rational $x$ with $|x|>0$, integration on each interval gives
$|{\cal F}f(x)|\le2S/(\sqrt{2\pi}|x|)$. Consequently
\begin{equation}\label{eq:fouriertail}
 \int_{|x|>R}|{\cal F}f(x)|^2\,dx\le {4S^2\over \pi R}.
\end{equation}
For example take integer $R\ge1+16S^2/\epsilon^2$; truncation then
costs at most $\epsilon/2$ in norm. On $[-R,R)$ use mesh and sample
errors making the local uniform error at most
$\epsilon/(4(R+1))$. Thus $\mathcal F f$ is a named vector, including
when the support contains $x=0$. The conjugate-sign transform
$\mathcal G$ has the same estimates.

\begin{theorem}[Plancherel for the step completion]\label{thm:pr-fourier}
$\mathcal F$ and $\mathcal G$ extend to mutually inverse PR isometries of
the represented Hilbert space. The conjugated Hamiltonian
$\mathcal F H_{p^2/2}\mathcal G$ acts as $-\frac12 d^2/dx^2$ on compact
piecewise-polynomial functions with continuous function and first
derivative, using their piecewise second derivative as the graph value
with the indicated factor.
\end{theorem}
\begin{proof}
We give a norm proof to establish the stability needed when passing
from finite input to an arbitrary Hilbert name. For a rational
$\alpha>0$, direct integration of the two half-lines gives
\[
 {1\over 2\pi}\int_{\mathbb R}e^{-\alpha|x|} e^{ixu}\,dx
       =K_\alpha(u):={1\over\pi}{\alpha\over \alpha^2+u^2}.
\]
The discarded tails on the left have absolute integral at most
$2e^{-\alpha R}/\alpha$ before the prefactor, and
$e^{-\alpha R}\le1/(1+\alpha R)$ suffices for a PR cutoff.
Multiplying the two copies of the finite-interval Fourier integral,
inserting this formula, and substituting $u=p-q$ gives
\begin{equation}\label{eq:abel-fourier}
 \int e^{-\alpha|x|}|{\cal F}f(x)|^2\,dx
     =\int K_\alpha(u)\langle f,T_u f\rangle\,du,\qquad
 (T_u f)(p)=f(p-u).
\end{equation}
Here are finite-sum details for these changes of variables. On
$|x|\le R$ and $|p|,|q|\le D$ containing the momentum support, split
the rectangle in $p,q$ at all step endpoints. The Lipschitz bounds of
$e^{-\alpha|x|}e^{ix(p-q)}$ in the three coordinates are
$\alpha+2D,R,R$. With maximal mesh $\eta$, the integrated error of
one rectangular sum is bounded by
$2R A^2(\alpha+2D+2R)\eta/(2\pi)$. Exchanging finite sums is exact.
Removing the $x$ cutoff costs at most $A^2 e^{-\alpha R}/(\pi\alpha)$.

For each pair of step intervals $I,J$, the substitution reduces to
\[
 \int_I\int_J K_\alpha(p-q)\,dq\,dp
 =\int K_\alpha(u)\,|I\cap(J+u)|\,du.
\]
Replace the kernel on $[-2D,2D]$ by rational mesh samples, to uniform
error $\eta/(\pi\alpha^2)$ plus sampling error, since
$|K_\alpha'|\le1/(\pi\alpha^2)$.
For one mesh interval $L$, the equality is equality of areas of rational
polygons. The linear map $(p,q)\mapsto(p,p-q)$ has determinant $-1$
and preserves triangle areas by the rational determinant formula. The
vertical slice length after the change is
$1_L(u)|I\cap(J+u)|$, a piecewise linear function whose integral is the
polygon area by summing trapezoids. The kernel-approximation error on
either side after summation is bounded by $A^2$ times its uniform error.
This proves \eqref{eq:abel-fourier} at arbitrary precision.

The kernel has integral one. Substitute $v=\alpha/u$ on $u\ge\alpha$.
Evenness gives the total integral as
$4\pi^{-1}\int_0^1(1+s^2)^{-1}\,ds=1$: integration of the finite
geometric identity leaves error at most $1/(2n+1)$ and yields the
alternating series defining $\pi/4$. Its mass on $|u|>\delta$
is at most $2\alpha/(\pi\delta)$.
Translations preserve the norm by translating the interval partition.
On each interval, its symmetric difference from its translate has length
at most $2|u|$, hence
$\|T_u f-f\|\le S\sqrt{2|u|}$. This estimate also constructs $T_u$
for real named $u$ by rational translations. It follows directly that
the right side of \eqref{eq:abel-fourier} differs from $\|f\|^2$ by
at most
\[
 \|f\|S\sqrt{2\delta}+{4\|f\|^2\alpha\over\pi\delta}.
\]
The left side differs from $\|{\cal F}f\|^2$ by at most
$\alpha R\,\|{\cal F}f\|^2+4S^2/(\pi R)$, using
$1-e^{-v}\le v$ for $v\ge0$. A bound for $\|{\cal F}f\|^2$ before
isometry is known is $(A^2+4S^2)/\pi$, by splitting at $R=1$.
For a requested error $0<\epsilon\le1$, take rational
$C\ge1+A+S+\|f\|$. One valid schedule is
\[
 R=\lceil32 C^2/\epsilon\rceil+1,\qquad
 \delta=\epsilon^2/(128 C^4),\qquad
 \alpha={\epsilon\delta\over64(R+1)C^2}.
\]
Indeed $\|{\cal F}f\|^2\le 5C^2$, and the four error terms above
sum to less than $\epsilon$. Each choice needs only rational
arithmetic, with optional bounded dyadic rounding
(for a rational $a/b>0$ with $a,b$ positive integers, $2^{-b}\le a/b$).
This proves $\|{\cal F}f\|=\|f\|$. The identical proof gives isometry
of $\mathcal G$.

An isometry sends a regular step name to a regular Hilbert name by
evaluating a more accurate input step and rounding the output within
the remaining budget. For compact steps $f,g$, integration over the
finite product of their supports gives
$\langle{\cal F}f,g\rangle=\langle f,{\cal G}g\rangle$ with our
inner-product convention. Extension by continuity gives the adjoint
identity everywhere. Preservation of inner products follows by
polarization from the norm identity. It implies
$\mathcal G\mathcal F=I$ and $\mathcal F\mathcal G=I$ by testing against all
step vectors. Finally two integrations by parts give
$\mathcal G(f'')(p)=-p^2\mathcal G f(p)$ for the stated compact
piecewise-polynomial $f$ (with named coefficients and certified
breakpoints). Interior boundary terms cancel by continuity of $f,f'$.
On each bounded momentum interval the transforms of $f,f''$ agree
with their integral formulas, by compact step approximation in both
$L^2$ and $L^1$ on the original bounded support. This proves the graph
equation, including the second component as a named Hilbert vector.
\end{proof}

Position in an interval $I$ now has momentum-space projection
$\mathcal G1_I\mathcal F$ and probability $\|1_I{\cal F}v\|^2$ in unit
input $v$. For $M_b f(p)=e^{ibp}f(p)$, direct substitution on steps
and contraction on all names give
\[
 T_aM_b=e^{-iab}M_bT_a,\qquad
 {\cal F}T_af(x)=e^{iax}{\cal F}f(x),\qquad
 {\cal F}M_bf(x)=({\cal F}f)(x+b).
\]
All identities are equalities of Hilbert names. The translation and
multiplication groups are strongly continuous by the preceding step
estimates. They therefore give the exponentiated canonical commutation
relations and both sharp position and momentum interval experiments
on a genuine continuous-spectrum system. In physical units the energy
phase is $e^{-itp^2/(2m\hbar)}$, the Fourier kernel is
$(2\pi\hbar)^{-1/2}e^{ixp/\hbar}$, and one supplies positive lower
bounds for $m,\hbar$.

\subsection{A particle in a bounded potential}

Let $H_0=\mathcal F H_{p^2/2}\mathcal G$ on position space. Supply a
bounded self-adjoint operator $V$ by its action on steps and an integer
bound $B\ge\|V\|$. A real bounded position potential with local PR
continuity moduli and rational sampling is such an input, as is a
bounded square well with rational endpoints by cutting intervals
exactly. Put $H=H_0+V$ on the graph domain of $H_0$.
Subtracting $Vx$ from the second component proves closedness, and the
adjoint identity reduces to that of $H_0$ by bounded self-adjointness
of $V$. Graph approximants to a finite step $f$ are explicit:
take $\mathcal F E_R\mathcal G f$, choosing $R$ from
\eqref{eq:fouriertail} for $\mathcal G f$. Its graph partner is
$\mathcal F((p^2/2) E_R\mathcal G f)+V\mathcal F E_R\mathcal G f$.
This gives exactly the density service in Section~\ref{sec:stone}.
We now construct the remaining resolvent data for unitary propagation.

Fix a rational $\rho>0$, either sign $\sigma\in\{-1,1\}$, and an
integer $L\ge\max(1,\rho,2B)$. In momentum coordinates
$R_0=(H_0+\sigma iL)^{-1}$ is the multiplier $(p^2/2+\sigma iL)^{-1}$,
transported by $\mathcal F$. It has bound $1/L$, and its graph
partner at input $f$ is $f-\sigma iL R_0 f$. The identity
\[
 H+\sigma iL=(I+V R_0)(H_0+\sigma iL)
\]
on graph inputs gives a two-sided inverse
$R=R_0\sum_{j\ge0}(-V R_0)^j$. The norm of $V R_0$ is at most $1/2$,
so the tail from $j=m$ is at most $2^{1-m}$ before the leading $R_0$.
The geometric identity for finite sums and this bound verify both inverse
equations. Every output is in the graph since $R_0$ is applied last.
Symmetry gives the sharper estimate $\|R\|\le1/L$ by
$\|(H+\sigma iL)v\|^2=\|Hv\|^2+L^2\|v\|^2$. Thus, with
$C=\sigma i(\rho-L)R$,
\[
 (H+\sigma i\rho)^{-1}
     = R\sum_{j\ge0}(-C)^j,\qquad \|C\|\le 1-\rho/L.
\]
With $b=\lceil L/\rho\rceil$, $(1-\rho/L)^b\le1/2$
since $(1-a)^b(1+ba)\le1$ for $0\le a\le1$, by finite induction.
The companion \texttt{ResolventTail.lean} proves its cleared-denominator
form $u^k(u+v+kv)\le(u+v)^{k+1}$ in the object calculus.
The tail from $j=bm$ is at most $(L/\rho)2^{-m}$, again before the
leading factor. The graph partner and both inverse identities follow
as above, and the inverse has bound $1/\rho$ by symmetry.
Products of a variable number of approximate operators here use fixed
per-factor precision, chosen by the finite telescoping estimate
following \eqref{eq:galerkin}; iteration on finite step codes is
ordinary primitive recursion. This supplies all imaginary-axis
resolvents, in particular at every positive integer scale, for the
resolvent construction of the unitary group. The Hamiltonian acts on the
compact polynomial test functions above by
$-\tfrac12 f''+Vf$, with the full Sobolev-type domain specified by
the transported multiplier graph.

\section{Identical particles}
\label{sec:q-identical}

Let the one-particle space have \(d\ge1\) specified orthonormal modes.
For \(N\) labelled particles a basis of the tensor power is indexed by
words \(w=(i_1,\ldots,i_N)\) in these modes; the empty word is the
vacuum. Permuting the slots by \(\pi\) defines a unitary \(U_\pi\),
with the convention \((\pi w)_s=w_{\pi^{-1}(s)}\).
For \(\chi=1\) or \(\chi=\operatorname{sgn}\), put
\[
               P_{N,\chi}=\frac1{N!}\sum_{\pi\in S_N}\chi(\pi)U_\pi .
\]
This rational matrix is a self-adjoint projection. Inversion preserves
\(\chi\), giving self-adjointness, and in the square every product
\(\tau\) occurs \(N!\) times with sign \(\chi(\tau)\).
Relabelling gives \(U_\tau P_{N,\chi}=\chi(\tau)P_{N,\chi}\).
Conversely, averaging fixes every vector with that transformation law.
The two choices specify the Bose and Fermi sectors, respectively; the
exchange law is a physical choice of sector. The finite construction may
test all of the \(N^N\) lists of length \(N>0\) for being a permutation
and count inversions for the sign. For \(N=0\), use one empty permutation.

\paragraph{Occupation coordinates.}
For bosons choose \(\alpha=(\alpha_1,\ldots,\alpha_d)\) in
\(\mathbb N^d\) with \(|\alpha|=\sum_j\alpha_j=N\).
There are \(m_\alpha=N!/\prod_j \alpha_j!\) words with these
multiplicities, by successive choices of slots. Their sum divided by
\(\sqrt{m_\alpha}\) is a unit vector \(|\alpha\rangle\).
These vectors are orthogonal because their supports are disjoint and
they span the symmetric sector because all coefficients on an orbit
are equal. For fermions a repeated letter forces the coefficient to
equal its negative under a transposition of equal slots. It therefore
vanishes. For each remaining increasing word take the signed sum of
its permutations divided by \(\sqrt{N!}\); on this orbit the
transformation law determines all coefficients from the increasing
word. This gives an orthonormal basis indexed by bit vectors
\(\alpha\in\{0,1\}^d\) of weight \(N\).
Consequently the sector dimensions are \({N+d-1\choose d-1}\) and
\({d\choose N}\), with the latter zero for \(N>d\). The first count
chooses the positions of \(d-1\) dividers in \(N+d-1\) slots; the
second chooses the occupied modes. Every denominator here has a positive
integer lower bound before taking its square root.

\paragraph{Creation, annihilation and domains.}
The Fock space is the Hilbert direct sum over \(N\) of these sectors.
Its effective vectors can be presented by finite occupation-coordinate
approximants with the usual \(\ell^2\) error witnesses. On finite
support set
\[
 \begin{aligned}
 a_j^*|\alpha\rangle&=\sqrt{\alpha_j+1}\,|\alpha+e_j\rangle,&
 a_j|\alpha\rangle&=\sqrt{\alpha_j}\,|\alpha-e_j\rangle
 &&\text{(bosons)} .
 \end{aligned}
\]
The term with \(\alpha_j=0\) in the annihilation formula is zero. For
fermions, creation is zero for an occupied mode and annihilation is zero
for an empty mode; otherwise change that bit and multiply by
\((-1)^{\sum_{k<j}\alpha_k}\). The two matrices are adjoints on the
finite-support domain by transposing each allowed basis transition.
For bosons these formulas give \([a_i,a_j^*]=\delta_{ij}I\) and
\([a_i,a_j]=[a_i^*,a_j^*]=0\). Indeed, on the same mode the two
number-preserving products have coefficients \(\alpha_i+1,\alpha_i\);
on different modes the two square-root factors are independent and the
same forbidden moves give zero in both orders. For fermions the
three formulas hold with anticommutators \(ST+TS\) in place of
commutators. On the same bit the two number-preserving products are its
empty and occupied projections and each doubled creation or annihilation
vanishes. On different bits the move at the smaller index flips the
sign for the move at the larger index while the reverse sign is unchanged.
An impossible move again gives zero in both orders.

For bosons the sector norms of creation and annihilation are bounded by
\(\sqrt{N+1}\) and \(\sqrt N\) because distinct admissible basis
vectors have distinct images and the squared coefficients satisfy these
bounds squared. On Fermi Fock space the corresponding bounds are one.
Thus an explicit sufficient domain witness for both Bose operations
is a cutoff \(M(p)\) with
\[
           \sum_{|\alpha|>M(p)} (|\alpha|+1)|c_\alpha|^2\le2^{-2p}.
\]
Each tail statement can equivalently be tested on every finite partial
sum. Apply the finite matrix to sectors through \(M(p+1)\) and allocate
the remaining half-error to the coefficients. For example take input
norm error at most \(2^{-p-1}/(M(p+1)+1)\) on this finite truncation.
The unboundedness is visible on states with \(N\) particles in mode
\(j\). A supplied one-particle unitary \(U\) acts on each sector by
\(U^{\otimes N}\), which commutes with permutations and is an isometry
with inverse \((U^*)^{\otimes N}\). It therefore acts on the direct sum
with the original tail bound.

For real mode energies \(\omega_j\), the free Hamiltonian on the
occupation basis multiplies by \(E_\alpha=\sum_j\omega_j\alpha_j\)
in units with \(\hbar=1\). It equals
\(\sum_j\omega_j a_j^*a_j\) on finite support, by the basis formulas.
Its unitary evolution multiplies each coefficient by
\(\exp(-itE_\alpha)\), preserving all unweighted tails and satisfying
the group law coordinatewise. For an effective derivative at time
\(t\), supply a cutoff for tails of
\(\sum_\alpha |E_\alpha c_\alpha|^2\). The tail error between the
difference quotient and \(-iE_\alpha\exp(-itE_\alpha)c_\alpha\) is at
most twice the norm of that weighted tail, since
\(|e^{ix}-1|\le |x|\). On a finite cutoff with \(|E_\alpha|\le K\),
the error is at most \(|h|K^2\|c\|_2/2\) for step \(h\), from
\(|e^{ix}-1-ix|\le x^2/2\) (integrate \(|e^{iy}-1|\le |y|\) on the
real segment). These are explicit input-dependent moduli for
Schr\"odinger's equation on the supplied domain.

\section{Open-system dynamics}
\label{sec:q-open}

For a finite-dimensional system an explicit time-independent
Markovian generator consists of a Hermitian matrix \(H\) and finitely
many matrices \(L_1,\ldots,L_r\). Write
\[
 \begin{aligned}
 M&=\sum_j L_j^*L_j, & \mathcal D(X)&=\sum_j L_jXL_j^*-\tfrac12(MX+XM),\\
 \mathcal H(X)&=-i(HX-XH), & \mathcal G&=\mathcal H+\mathcal D.
 \end{aligned}
\]
The following argument identifies the whole finite-dimensional class
when an effective right derivative at zero is part of the semigroup
data. For superoperators use the estimate
\(\|\mathcal A(X)\|\le b\|X\|\) with matrix operator norm estimates;
every use can equivalently be read as a bounded-vector test.

\begin{theorem}[Finite Lindblad evolution]\label{thm:q-lindblad}
For \(t\ge0\), \(S_t=\exp(t\mathcal G)\) is a channel,
\(S_{t+s}=S_t S_s\), and \(d S_t/dt=\mathcal G S_t\) (from the right
at zero). On a bounded interval it admits a primitive-recursive
approximation by channels that preserve trace at every step.
\end{theorem}
\begin{proof}
Choose integers \(a\ge2\|H\|\), \(\ell\ge\sum_j\|L_j\|^2\),
\(C=1+a+2\ell\), and \(T\ge\max(1,t)\).
For \(h=t/n\), \(n\ge CT\), take
\[
 K_0=(I-hM)^{1/2},\quad K_j=\sqrt h\,L_j,\quad
 U_h=\exp(-ihH),\qquad
 F_h(X)=U_h\left(\sum_{j=0}^{r}K_j X K_j^*\right)U_h^* .
\]
The input \(I-hM\) is positive since \(0\le M\le\ell I\). These
are Kraus matrices followed by a unitary, and
\(\sum_{j=0}^r K_j^*K_j=I\), so every \(F_h\) is a channel.
The binomial series of Theorem~\ref{thm:q-sqrt} gives
\[
 \|K_0-I\|\le h\ell,\qquad
 \|K_0-I+hM/2\|\le h^2\ell^2 .
\]
Indeed use \(F_m(hM)\) from that proof, whose nonconstant absolute
coefficients sum to at most one and whose linear coefficient is
\(-1/2\). Expanding \(K_0XK_0\) about \(X\) shows that the channel
inside parentheses differs from \(I+h\mathcal D\) by at most
\(3h^2\ell^2\) in superoperator norm. Here and below the identity
superoperator is denoted by \(I\) when it acts on matrices.

Unitary conjugation \(W_h(X)=U_hXU_h^*\) equals \(e^{h\mathcal H}(X)\):
expand the two exponentials and use the binomial formula for the
commuting operations of left and right multiplication. The same
identity gives unitarity by multiplication of the one-matrix series.
For \(0\le hC\le1\), the factorial series gives
\(\|W_h-I\|\le3ha\) and
\(\|W_h-I-h\mathcal H\|\le2h^2a^2\).
As \(\|\mathcal D\|\le2\ell\), multiplying the two approximations
bounds \(\|F_h-I-h\mathcal G\|\) by
\(3h^2\ell^2+2h^2a^2+6h^2a\ell\le 8h^2 C^2\).
The factorial tail of \(e^{h\mathcal G}\) beyond its linear part is
at most \(2h^2 C^2\). Also \(\|F_h\|\le1+h\ell\le e^{hC}\), since
\(K_0\) is a contraction, and \(\|e^{h\mathcal G}\|\le e^{hC}\).
Telescoping the difference of the \(n\)-th powers yields
\[
             \|F_{t/n}^{\,n}-e^{t\mathcal G}\|
               \le 10 C^2T^2\,4^{CT}/n .
\]
The factor \(4^{CT}\) bounds \(e^{CT}\) since \(C,T\) are integers.
Taking \(n\ge CT+10C^2T^2 4^{CT}2^{p+1}\) leaves half the precision
budget for finite matrix arithmetic. The limit has positive Choi
matrix by entry convergence on the finitely many matrix units, and
preserves trace by the same test. Thus it is a channel by
Theorem~\ref{thm:q-choi}.

For explicit series control at \(\|A\|\le D\), \(D\ge1\) integer,
let \(k=2D\). Beyond term \(k\) consecutive absolute terms have ratio
at most \(1/2\); term \(k\) is at most \(4^D\) by the scalar
exponential bound. Truncating after index \(k+p+2D+2\) therefore
makes the tail less than \(2^{-p}\). This estimate justifies
Cauchy products by finite truncation. The binomial identity for
commuting exponents gives the semigroup law; the linear Taylor
remainder bound above, with \(\mathcal G\) in place of \(\mathcal H\),
gives the derivative at each time after multiplication by \(S_t\).
\end{proof}

\begin{theorem}[Recovery of the generator]\label{thm:q-lindblad-converse}
Suppose \(T_0=I\) and \(T_t\) is a family of channels in dimension \(d\)
with a named right derivative \(G\) at zero and a primitive-recursive
accuracy modulus on each matrix unit. Then \(G\) has the above Lindblad
form with at most \(d^2\) operators \(L_j\). If the family is a semigroup,
\(T_t=e^{tG}\).
\end{theorem}
\begin{proof}
For a linear map \(A\) use its Choi matrix \(J_A\) from
Theorem~\ref{thm:q-choi}. Set \(w=\sum_i e_i\otimes e_i\),
\(P=ww^*/d\), \(Q=I-P\). Since \(J_I=ww^*\) and \(Qw=0\),
\[
       Q J_G Q=\lim_{h\downarrow0} Q J_{T_h}Q/h\ \ge0 .
\]
Take the limit along positive dyadics with the given modulus.
All \(J_{T_h}\) are Hermitian, so \(J_G\) is Hermitian; also
\(\operatorname{tr}G(X)=0\) by differentiating trace preservation.
Factor the positive matrix \(QJ_GQ\) with its positive square root
and reshape its \(d^2\) columns into matrices \(L_j\), as in the
Choi proof. Their map has Choi matrix \(QJ_GQ\).
Set
\[
             v=J_Gw/d-\frac{w^*J_Gw}{2d^2}w,
             \qquad V_{ai}=v_{ia}.
\]
Multiplication gives
\(J_G-QJ_GQ=J_GP+PJ_G-PJ_GP=vw^*+wv^*\). Testing matrix units,
the latter is the Choi matrix of \(X\mapsto VX+XV^*\).
Thus \(G(X)=\sum_j L_jXL_j^*+VX+XV^*\). Its vanishing trace on
matrix units gives \(V+V^*=-M\). With
\(H=(V^*-V)/(2i)\) this is the displayed generator.

For uniqueness from the semigroup law one can use finite telescoping.
In dimension \(d\), estimates on all units within \(\delta/d^2\)
give a superoperator bound \(\delta\), since
\(|X_{ij}|\le\|X\|\). Hence for each rational \(0<\delta\le1\) the
derivative modulus and the exponential Taylor bound give a positive
dyadic cutoff on \(h\) with
\(\|T_h-e^{hG}\|\le h\delta\) and
\(\|T_h\|\le 1+h(\|G\|+1)\) (use the derivative estimate directly for
the latter). These two closed inequalities also hold when \(h=0\)
by \(T_0=I\). For a general nonnegative named \(h\) below cutoff, a
strict violation of a tested inequality would exclude \(h>0\).
Since \(h\ge0\), this implies \(h=0\), contradicting the same violation.
This establishes the closed inequalities without deciding equality
to zero. For \(0\le t\le T\), take \(n\) large enough that \(h=t/n\)
lies under the cutoff. The difference of powers is bounded by
\(T e^{T(b+1)}\delta\) with an integer \(b\ge\|G\|\). Thus
\(T_t=e^{tG}\), with the required \(n\) obtained from the supplied
derivative modulus and the bound on \(T\). The modulus here applies to
every sufficiently small positive time, as in the definition of a
right derivative.
\end{proof}

\begin{example}[Spontaneous decay of one excited level]
In the ordered basis \(g,e\), take \(H=\omega |e\rangle\langle e|\)
and \(L=\sqrt\gamma\,|g\rangle\langle e|\) with \(\gamma\ge0\).
Set \(z=e^{-(\gamma/2+i\omega)t}\), \(q=|z|^2=e^{-\gamma t}\).
The evolved density matrix is
\[
 \begin{pmatrix}
 \rho_{gg}+(1-q)\rho_{ee} & \overline z\,\rho_{ge}\\
 z\rho_{eg} & q\rho_{ee}
 \end{pmatrix}.
\]
It follows either by multiplication with Kraus matrices
\(K_0=|g\rangle\langle g|+z|e\rangle\langle e|\) and
\(K_1=\sqrt{1-q}\,|g\rangle\langle e|\), or by differentiating:
the excited population solves \(x'=-\gamma x\), the ground population
gains \(\gamma x\), and the \(ge\) entry has coefficient
\(i\omega-\gamma/2\). The formula composes in time by multiplication
of \(z\) and \(q\), and agrees with the identity at zero. It includes
vanishing decay rate and vanishing initial excited population with the
same formulas.
\end{example}

\section{Formalization, numerical evidence and remaining scope}\label{sec:lean}
\subsection{What the Lean kernel checks in this packet}
The proof target is ProveIt's ordinary judgment $\Gamma\vdash_{\mathrm{PRAAx}}\varphi$, represented in Lean by \texttt{Prf PRAAx}. The repository supplies a kernel-checked arithmetic library and a separate collection of research reports with varying verification status. We use the former for the endpoints below. The analytic arguments in Sections~\ref{sec:finite-quantum}--\ref{sec:q-open} are mathematical proofs with explicit realization data; only the following individually named results are supplied as closed derivations. File and module names for replay are in the archive.

\begin{center}\small\renewcommand{\arraystretch}{1.22}
\begin{tabularx}{\linewidth}{@{}p{.29\linewidth}X@{}}\toprule
Closed Lean theorem & First-order mathematical statement \\ \midrule
\texttt{ports\_closed} & The sum of the two unnormalized output weights $|z+w|^2+|z-w|^2$ is twice the input weight.\\
\texttt{spin\_closed} & The map $(z,w)\mapsto(az-\bar b w,bz+\bar a w)$ scales squared norm by $|a|^2+|b|^2$.\\
\texttt{lagrange\_closed} & $|\bar a z+\bar b w|^2+|aw-bz|^2=(|a|^2+|b|^2)(|z|^2+|w|^2)$.\\
\texttt{stability\_closed} & If $a,b,c,d$ are signed integer components of size at most $M$, with $|a-c|,|b-d|\le e$, the two squared complex amplitudes differ by at most $4Me$.\\
\texttt{closed\_bound} & $a^k(a+b+kb)\le(a+b)^{k+1}$. This is $(1-r)^k(1+kr)\le1$ for $r=b/(a+b)$ after clearing a positive denominator.\\ \bottomrule
\end{tabularx}\end{center}
Each statement universally quantifies its natural coordinates. A signed component is a pair of terms and its equality is cross-addition of their parts. Taking common positive denominators in the first four equations/inequalities recovers the rational uses. The first three proofs turn an integer polynomial identity into an object derivation via \texttt{ipra\_ring}. The fourth constructs nonnegative slack terms using object order. The fifth carries out object-theoretic quantifier-free induction on $k$. The kernel's axiom reports for these endpoints list \texttt{propext}, \texttt{Classical.choice} and \texttt{Quot.sound} in Lean's metatheory. Their conclusions are derivations in the intuitionistic object theory. No analytic theorem is imported into it from a classical Hilbert-space library.

\paragraph{An explicit object-language interface.}
\texttt{QuantumNames.lean} defines \texttt{VectorCode} by six actual PR function symbols. A length and a denominator predecessor are functions of the stage; four coordinate readers give positive and negative parts for real and imaginary numerators as functions of index and stage. Write $q_s=D_s^{-1}(u_{is}+iv_{is})_i$, zero-padding after its length, where $D_s\ge1$. If $q_s,r_t$ have lengths $l,l'$ and denominators $D,F$, then the literal meaning of the compiled comparison \texttt{within} with radius $p/b$ ($b>0$) is
\[
 b^2\sum_{i<l+l'}\!\left[(Fu_i-Du'_i)^2+(Fv_i-Dv'_i)^2\right]\ \le\ (pDF)^2.
\]
Each signed difference is calculated using the sum of the two opposite truncated subtractions of its unsigned sides. A zero-one factor enforces the length independently for each reader, so that surplus reader values never enter the distance. The sum is an actual primitive recursion on its length with two stage parameters. The Lean definitions \texttt{valid} and \texttt{equivalent} close the norm bound, the regularity test with $b=2^s$, or the comparison of equal-stage samples with $p=2$, respectively. \texttt{within\_qf} checks that the test is quantifier-free. The structure \texttt{CertifiedVector} requires an actual derivation of validity. These definitions provide syntax; the constructor library in the next subsection remains to be implemented.

\subsection{A path to full object-language proofs}
The interface just given handles fixed symbols. Uniform families add ordinary numerical arguments, or pair the tuple of arguments into one extra parameter. A column sampler has arguments $(j,s)$ before its coordinate index. An operator action on a finite rational vector has arguments $(q,k)$ and returns a finite vector code. A graph approximation returns a pair of names through jointly PR samplers indexed by the input and target precision. The following develops the proof modules in dependency order. Each obligation has a finite assertion or an explicit instantiation rule against which an implementation can be tested.

\begin{enumerate}
\item \textbf{Rationals and finite arrays.} Implement the codes of Section~\ref{sec:arithmetic}, list projections, zero-padding, common-denominator products, matrix products, contractions, partial sums of squared absolute values and bisection. For the primitive-recursive loops use \texttt{recDef}/\texttt{prec}, the base and successor equations and \texttt{Core.Order}. Derive the read-after-write equations and the fold invariant for each accumulator. Prove the general finite Lagrange identity by adding one coordinate and expanding the new cross terms. The signed ring backend proves each local identity; \texttt{lagrange\_closed} is its two-coordinate test. The order library converts sums of squares to nonnegative terms. This gives squared triangle tests, avoiding a square-root comparison in the object formula.
\item \textbf{Hilbert setoids and operator bounds.} For every constructor produce its sampler and a proof from the input \texttt{valid} assumptions to output validity. Use the displayed stages in Section~\ref{sec:infinitehilbert}. For an error estimate on the limit, fix an output stage and derive the rational comparison on two approximants. The closed-inequality rule uses a user-independent PR function of their denominators and the claimed error constant. For \eqref{eq:H1} introduce one finite input and one observation precision. Prove extension, adjointness and extensionality by two column truncations. The rate for a series becomes the PR cutoff $T(k)$, and the trace equality compares two finite basis traces at the maximum of their cutoffs. The output of completed constructions remains an ordinary sampler. Formalize the diagonal lemma once and instantiate it for each of these spaces, including compactly supported step functions.
\item \textbf{States and outcome probabilities.} Encode a density by its finite entries or a purified array and its normalization and positivity tests. For a given finite algebraic operation, clear denominators and induct over the matrix contraction index. For variable-size instruments prove the partial Kraus-sum identity before taking the full length. For conditioning add the explicit lower-bound test on the branch weight. For Choi factorization implement the square-root coefficients via their binomial recurrence, prove the coefficient inequalities and residual polynomial identity, then invoke the operator completion rule with $m=L^2 2^{2k+2}$. The indices $(i,a)$ and the column index of the root are the Kraus matrix indices, so the proof generator never needs to branch on rank. To mechanize effect representation, instantiate the supplied evaluator on the finite shifted matrix units. For the repeat-addition argument package the $n$ applications on $A/(nm)$ into one jointly sampled family, allocate $\delta_k/(1+mn)$ to each use, and induct over its bounded prefix at that precision.
\item \textbf{Continuous and sharp observables.} For the tables $W_{Nj}$, prove Pascal's recurrence, normalization, and the two factorial moments by finite bounded sums. The coupling estimate in Lemma~\ref{lem:binomial} may be formalized by counting success tuples in $\{0,\ldots,N-1\}^a$ and bounding collision tuples by ${a\choose2}N^{a-1}$ for $a\ge1$; treat $a=0$ separately. Multiply probabilities by their common finite denominator. The Chebyshev bound is termwise comparison of weighted squares. Feed the two resulting approximation inequalities to \texttt{within} on a finite basis test. For a sharp cut add either the quadratic-form certificate $B^2\ge g^2 I$ or a scalar boundary-mass modulus; use \texttt{closed\_bound} to eliminate a geometric tail. For unbounded $H$ implement the strong sampler for $F$ by dense range witnesses, and compare $g_a(F)$ on $Dy$ as in Theorem~\ref{thm:unbounded-calculus}. No spectral enumeration occurs.
\item \textbf{Generators and integrals.} Fix PR samplers for both resolvents indexed by $R$ and for the graph-density approximant. The graph at $(x,y)$ can be tested by $R_+(1)(y+ix)=x$. The analytic assertions reduce to tests at a finite input vector and precision; the resolvent identities are derived from the inverse equations on range witnesses. The key convergence loop computes $v=D_bx$, its second graph bound, the cutoff $R$, and the degrees of the two exponential approximations, in that order. Realize a variable power by iteration on finite vectors with a fixed per-step budget (Section~\ref{sec:arithmetic}). For the inverse theorem implement vector-valued rational Riemann sums. A binary common grid proves the Cauchy comparison; the tail and mesh in the Laplace integral are functions of $R,k$ and the orbit modulus. Derive the translation identity first on a finite interval, with two endpoint strips bounded by their lengths times a uniform bound. The derivative formula is a closed remainder inequality $\|U(h)x-x+ihHx\|\le|h|\delta_k$ under $|h|\le2^{-d(k)}$; it contains no division by a potentially zero named increment. Implement the driven-system mesh comparison by the unitary telescope on a common refinement.
\item \textbf{Models.} Encode a step by its sorted rational endpoints and Gaussian-rational heights. Refinement preserves its rational inner products exactly. For Plancherel, implement in order the tail bound \eqref{eq:fouriertail}, the rectangular integral identity \eqref{eq:abel-fourier}, the rational polygon shear, and the four-term Abel error estimate at its displayed $(R,\delta,\alpha)$. For banded dynamics implement the stopped-word count as a bounded subset enumeration and \eqref{eq:galerkin} as a pairwise inequality on two finite cutoffs. In occupation coordinates prove the creation/annihilation relations on a single basis word before extending to the invariant finite core. For an open system prove the one-step $O(h^2)$ Kraus error, then the $n$-fold telescope of Theorem~\ref{thm:q-lindblad}. In the converse, differentiability is a supplied finite list of common-modulus estimates on matrix units. Each family can then be closed over its numerical parameters using \texttt{closeUniv}. The optional translation to the repository's native scalar calculus uses \texttt{pra\_iff\_native} only after obtaining the ordinary derivation.
\end{enumerate}

\subsection{A reproducible infinite-system prediction}
For the anharmonic ladder of Section~\ref{sec:dynamics} start in level $0$, put $t=1/2$, $g=1/4$, and measure the projection onto level $0$. The exact rational program \texttt{band\_evolution.py} gives the outward-rounded enclosure
\[
  0.98571462780342863\ \le\ |\langle e_0,U(1/2)e_0\rangle|^2\ \le\ 0.98572609533871460 .
\]
Seven levels and Taylor degree 67 suffice. Here $B=2|g|=1/2$, $m=6$, and the discarded-level error in the amplitude is at most $2(B|t|)^6/6!$. The factor $|t|$ in this sharper use of \eqref{eq:galerkin} follows directly from the stopped-word count, which has $(B|t|)^m$ before taking the time supremum. On the seven levels use $a=|t|((6+1/2)^2+B)$. The Taylor tail beginning at $r=68$ is at most $2a^{68}/68!$ since subsequent ratios are below $1/2$. The program multiplies the tridiagonal rational matrix into the previous vector, scales by $t/r$, and accumulates its zeroth coordinate with phase $(-i)^r$. If $\eta$ is the sum of the two errors then the squared modulus differs from that of the computed amplitude by at most $(2+\eta)\eta$ by unitarity. Thus the stated bounds control every discarded level of the infinite system. The companion \texttt{examples.py} checks Bell conditioning, the bias of a finite Bernstein measurement, and the boundary term $I-nP_{\rm top}$ in a finite oscillator commutator.

\subsection{Scope and further research}
The finite-dimensional theory allows arbitrary named complex entries and arbitrary finite dimensions. Separable systems require their basis or an equivalent dense presentation. Normal states, represented observables, continuous energy responses and unitary groups come with the constructive data the proofs use. The distinction is experimentally meaningful: continuous detectors have predictions for every presented state; a sharp spectral question can decide halting even for a scalar Hamiltonian (Proposition~\ref{prop:cut}). Likewise the oscillator has a strong orbit modulus depending on the state; no single time modulus controls the displacement of the whole unit ball for arbitrarily high energy. For example on the $j$th mode of the number operator the phase at $t=\pi/j$ is $-1$.

The theory here concerns nonrelativistic quantum mechanics. Interesting extensions are a form-domain construction of resolvents for singular Schr\"odinger potentials, scattering with PR propagation and wave-operator moduli, and unbounded Markov generators with an effective non-explosion criterion. These need estimates beyond the bounded potential and finite Lindblad results above. A more immediate formalization target is the operator-valued Bernstein norm inequality: once its finite counting proof and the Hilbert completion lemmas become \texttt{Prf PRAAx} endpoints, a single kernel-checked interface would serve square roots, continuous detector responses and energy cutoffs.

\begin{thebibliography}{9}
\bibitem{proveit} V. Reshetnikov et al., \emph{ProveIt}, Logic/PrimitiveRecursiveArithmetic. Ordinary iPRA is defined in \texttt{ConciseIPRA.PR}; the integer, order, power and bounded-sum modules are in \texttt{IPRADerivations.Core}. Repository snapshot \texttt{84192061db7eb09e923bd4ded5199f23d30c6e83}. \href{https://github.com/VladimirReshetnikov/ProveIt/tree/84192061db7eb09e923bd4ded5199f23d30c6e83/Logic/PrimitiveRecursiveArithmetic}{Pinned source and verification-status guide}.
\bibitem{busch} P. Busch, Quantum states and generalized observables: a simple proof of Gleason's theorem. \emph{Phys. Rev. Lett.} \textbf{91} (2003), 120403. Classical context for the all-effects additivity assumption in Theorem~\ref{thm:q-effect}. \href{https://doi.org/10.1103/PhysRevLett.91.120403}{doi:10.1103/PhysRevLett.91.120403}.
\bibitem{watrous} J. Watrous, \emph{The Theory of Quantum Information}, Cambridge University Press, 2018, chapters 2--3. Finite channel and discrimination terminology. \href{https://cs.uwaterloo.ca/~watrous/TQI/}{Author's text and errata}.
\bibitem{lindblad} G. Lindblad, On the generators of quantum dynamical semigroups. \emph{Comm. Math. Phys.} \textbf{48} (1976), 119--130. Classical generator theorem underlying the finite-dimensional formulation of Section~\ref{sec:q-open}. \href{https://doi.org/10.1007/BF01608499}{doi:10.1007/BF01608499}.
\bibitem{brattka} V. Brattka and R. Dillhage, Computability of the spectrum of self-adjoint operators. \emph{J. Universal Computer Science} \textbf{11} (2005), 1884--1900. Computability of spectral information for different operator presentations. \href{https://doi.org/10.3217/jucs-011-12-1884}{doi:10.3217/jucs-011-12-1884}.
\bibitem{ye} F. Ye, Toward a constructive theory of unbounded linear operators. \emph{J. Symbolic Logic} \textbf{65} (2000), 357--370. Constructive spectral and Stone theorems in Bishop's setting. \href{https://doi.org/10.2307/2586543}{doi:10.2307/2586543}.
\bibitem{hellman} G. Hellman, Constructive mathematics and quantum mechanics: unbounded operators and the spectral theorem, in \emph{Mathematics and Its Logics}, Cambridge University Press, 2021, pp. 212--236. Background on the constructive representation problem. \href{https://doi.org/10.1017/9781108657419.014}{doi:10.1017/9781108657419.014}.
\end{thebibliography}

\end{document}
